\documentclass[11pt]{article}

% ---- theme variables (passed via pandoc -V) ----
\newcommand{\ThemeName}{Concurrency \& Multithreading}
\newcommand{\ThemeKicker}{Systems Track}
\newcommand{\ThemeNumber}{01}

\usepackage{interviewstyle}
\setthemecolor{6366f1}{8b5cf6}

% pandoc-required plumbing
\providecommand{\tightlist}{\setlength{\itemsep}{1pt}\setlength{\parskip}{0pt}}
\setlength{\emergencystretch}{3em}
\providecommand{\pandocbounded}[1]{#1}

% syntax-highlighting macros emitted by pandoc (skylighting)
\usepackage{color}
\usepackage{fancyvrb}
\newcommand{\VerbBar}{|}
\newcommand{\VERB}{\Verb[commandchars=\\\{\}]}
\DefineVerbatimEnvironment{Highlighting}{Verbatim}{commandchars=\\\{\}}
% Add ',fontsize=\small' for more characters per line
\usepackage{framed}
\definecolor{shadecolor}{RGB}{35,38,41}
\newenvironment{Shaded}{\begin{snugshade}}{\end{snugshade}}
\newcommand{\AlertTok}[1]{\textcolor[rgb]{0.58,0.85,0.30}{\textbf{\colorbox[rgb]{0.30,0.12,0.14}{#1}}}}
\newcommand{\AnnotationTok}[1]{\textcolor[rgb]{0.25,0.50,0.35}{#1}}
\newcommand{\AttributeTok}[1]{\textcolor[rgb]{0.16,0.50,0.73}{#1}}
\newcommand{\BaseNTok}[1]{\textcolor[rgb]{0.96,0.45,0.00}{#1}}
\newcommand{\BuiltInTok}[1]{\textcolor[rgb]{0.50,0.55,0.55}{#1}}
\newcommand{\CharTok}[1]{\textcolor[rgb]{0.24,0.68,0.91}{#1}}
\newcommand{\CommentTok}[1]{\textcolor[rgb]{0.48,0.49,0.49}{#1}}
\newcommand{\CommentVarTok}[1]{\textcolor[rgb]{0.50,0.55,0.55}{#1}}
\newcommand{\ConstantTok}[1]{\textcolor[rgb]{0.15,0.68,0.68}{\textbf{#1}}}
\newcommand{\ControlFlowTok}[1]{\textcolor[rgb]{0.99,0.74,0.29}{\textbf{#1}}}
\newcommand{\DataTypeTok}[1]{\textcolor[rgb]{0.16,0.50,0.73}{#1}}
\newcommand{\DecValTok}[1]{\textcolor[rgb]{0.96,0.45,0.00}{#1}}
\newcommand{\DocumentationTok}[1]{\textcolor[rgb]{0.64,0.20,0.25}{#1}}
\newcommand{\ErrorTok}[1]{\textcolor[rgb]{0.85,0.27,0.33}{\underline{#1}}}
\newcommand{\ExtensionTok}[1]{\textcolor[rgb]{0.00,0.60,1.00}{\textbf{#1}}}
\newcommand{\FloatTok}[1]{\textcolor[rgb]{0.96,0.45,0.00}{#1}}
\newcommand{\FunctionTok}[1]{\textcolor[rgb]{0.56,0.27,0.68}{#1}}
\newcommand{\ImportTok}[1]{\textcolor[rgb]{0.15,0.68,0.38}{#1}}
\newcommand{\InformationTok}[1]{\textcolor[rgb]{0.77,0.36,0.00}{#1}}
\newcommand{\KeywordTok}[1]{\textcolor[rgb]{0.81,0.81,0.76}{\textbf{#1}}}
\newcommand{\NormalTok}[1]{\textcolor[rgb]{0.81,0.81,0.76}{#1}}
\newcommand{\OperatorTok}[1]{\textcolor[rgb]{0.81,0.81,0.76}{#1}}
\newcommand{\OtherTok}[1]{\textcolor[rgb]{0.15,0.68,0.38}{#1}}
\newcommand{\PreprocessorTok}[1]{\textcolor[rgb]{0.15,0.68,0.38}{#1}}
\newcommand{\RegionMarkerTok}[1]{\textcolor[rgb]{0.16,0.50,0.73}{\colorbox[rgb]{0.08,0.19,0.26}{#1}}}
\newcommand{\SpecialCharTok}[1]{\textcolor[rgb]{0.24,0.68,0.91}{#1}}
\newcommand{\SpecialStringTok}[1]{\textcolor[rgb]{0.85,0.27,0.33}{#1}}
\newcommand{\StringTok}[1]{\textcolor[rgb]{0.96,0.31,0.31}{#1}}
\newcommand{\VariableTok}[1]{\textcolor[rgb]{0.15,0.68,0.68}{#1}}
\newcommand{\VerbatimStringTok}[1]{\textcolor[rgb]{0.85,0.27,0.33}{#1}}
\newcommand{\WarningTok}[1]{\textcolor[rgb]{0.85,0.27,0.33}{#1}}

\begin{document}

\makecoverpage

\thispagestyle{empty}
{\hypersetup{hidelinks}\tableofcontents}
\clearpage

\begin{quote}
\textbf{How to use this file:} Read top-to-bottom for deep understanding. Jump to \hyperlink{section.9}{\textcolor{acc}{§9}} (Interview Questions) + \hyperlink{section.14}{\textcolor{acc}{§14}} (Cheat Sheet) for last-minute revision.
\end{quote}

\section{Overview --- What it is}\label{overview--what-it-is}

\textbf{Concurrency} is the ability of a system to deal with multiple tasks at the same time --- tasks make \emph{progress} simultaneously, but not necessarily \emph{execute} at the same instant.

\textbf{Parallelism} is actual simultaneous execution on multiple CPU cores.

\setcodelang{}

\begin{verbatim}
Concurrency (1 core):         Parallelism (2 cores):

Core 0: [T1][T2][T1][T2]      Core 0: [T1][T1][T1]
         ↑ interleaved          Core 1: [T2][T2][T2]
                                        ↑ truly simultaneous
\end{verbatim}

\begin{longtable}[]{@{}
  >{\raggedright\arraybackslash}p{(\columnwidth - 4\tabcolsep) * \real{0.1471}}
  >{\raggedright\arraybackslash}p{(\columnwidth - 4\tabcolsep) * \real{0.3980}}
  >{\raggedright\arraybackslash}p{(\columnwidth - 4\tabcolsep) * \real{0.4549}}@{}}
\toprule\noalign{}
\rowcolor{acc!16}
\begin{minipage}[b]{\linewidth}\raggedright
Dimension
\end{minipage} & \begin{minipage}[b]{\linewidth}\raggedright
Concurrency
\end{minipage} & \begin{minipage}[b]{\linewidth}\raggedright
Parallelism
\end{minipage} \\
\midrule\noalign{}
\endhead
\bottomrule\noalign{}
\endlastfoot
Goal & Structure / responsiveness & Speed / throughput \\
\rowcolor{zebra}
Requires & Scheduler, context switching & Multiple CPU cores \\
Example & Single-threaded event loop & Matrix multiplication on 8 cores \\
\rowcolor{zebra}
Risk & Race conditions, deadlocks & False sharing, synchronization \\
\end{longtable}

\textbf{Key mental model:} Concurrency is about \emph{design}; parallelism is about \emph{execution}. A single-core machine can be concurrent but never parallel.

A \textbf{thread} is the smallest unit of CPU scheduling --- it shares the process\textquotesingle s memory space (heap, globals) but has its own stack and program counter.

A \textbf{process} is an isolated address space; OS-level container for threads.

\setcodelang{}

\begin{verbatim}
Process
├── Shared: Heap, Code, File Descriptors, Global Variables
├── Thread 1: Stack, Registers, PC
├── Thread 2: Stack, Registers, PC
└── Thread 3: Stack, Registers, PC
\end{verbatim}

\section{Why It Exists}\label{why-it-exists}

\subsection{The Hardware Reality}\label{the-hardware-reality}

Modern CPUs stopped getting faster clock speeds around 2004 (the "Power Wall"). Instead, they grew more cores. A program that uses only one thread leaves 7 of 8 cores idle. Concurrency is the answer to: \emph{how do we exploit all this hardware?}

\subsection{The Latency Reality}\label{the-latency-reality}

I/O (disk, network, DB) is 10,000--1,000,000x slower than CPU. A server that waits synchronously for each DB query handles 1 req/s instead of 10,000 req/s. Concurrency lets the CPU work on other requests while one is waiting on I/O.

\setcodelang{}

\begin{verbatim}
Without concurrency:
Request 1: [CPU work]---[WAIT for DB]---[CPU work]
Request 2:                               [CPU work]---[WAIT for DB]...

With concurrency:
Request 1: [CPU work]---[WAIT]---[CPU work]
Request 2:          [CPU work]---[WAIT]---[CPU work]
                    ↑ overlap here — same total time, 2x throughput
\end{verbatim}

\subsection{Timeline of Need}\label{timeline-of-need}

\begin{itemize}
\tightlist
\item
  \textbf{1960s}: Batch OS needed to overlap I/O and computation
\item
  \textbf{1990s}: Web servers needed to handle many simultaneous users
\item
  \textbf{2004}: Multi-core CPUs made parallelism mainstream
\item
  \textbf{2010s}: Microservices, distributed systems, reactive programming
\item
  \textbf{Today}: GPUs, async runtimes, lock-free data structures
\end{itemize}

\section{Why FAANG Cares}\label{why-faang-cares}

\begin{longtable}[]{@{}
  >{\raggedright\arraybackslash}p{(\columnwidth - 2\tabcolsep) * \real{0.1188}}
  >{\raggedright\arraybackslash}p{(\columnwidth - 2\tabcolsep) * \real{0.8812}}@{}}
\toprule\noalign{}
\rowcolor{acc!16}
\begin{minipage}[b]{\linewidth}\raggedright
Company
\end{minipage} & \begin{minipage}[b]{\linewidth}\raggedright
Where Concurrency Is Critical
\end{minipage} \\
\midrule\noalign{}
\endhead
\bottomrule\noalign{}
\endlastfoot
\textbf{Google} & Search index serving (billions of queries/day), Bigtable/Spanner distributed transactions, MapReduce parallelism \\
\rowcolor{zebra}
\textbf{Meta} & Thrift RPC servers handle millions of concurrent connections; News Feed ranking runs parallel ML inference pipelines \\
\textbf{Amazon} & DynamoDB uses lock-free skip lists; Lambda cold starts require efficient thread/process pool management \\
\rowcolor{zebra}
\textbf{Apple} & Grand Central Dispatch (GCD) --- their concurrency framework; Core Data multi-context concurrency \\
\textbf{Microsoft} & Azure Service Bus, .NET async/await model, SQL Server lock manager, Windows thread scheduler \\
\rowcolor{zebra}
\textbf{Uber} & Real-time dispatch: matching riders/drivers under microsecond SLAs; surge pricing uses atomic counters \\
\textbf{Databricks} & Spark executor thread pools, Delta Lake transaction log with optimistic concurrency, shuffle service parallelism \\
\end{longtable}

\textbf{Interviewers test this because:}

\begin{enumerate}
\def\labelenumi{\arabic{enumi}.}
\tightlist
\item
  Concurrency bugs are subtle, hard to reproduce, and catastrophic in production
\item
  Senior engineers must design concurrent systems correctly from the start
\item
  It reveals depth of understanding of OS, memory models, and hardware
\end{enumerate}

\section{Core Concepts}\label{core-concepts}

\subsection{Threads vs Processes}\label{threads-vs-processes}

\textbf{Process:} OS-level isolation. Memory, file handles, and address space are separate. Context switch is expensive (\textasciitilde1--10 µs, requires kernel mode switch, TLB flush).

\textbf{Thread:} Lightweight. Shares process memory. Context switch is cheaper (\textasciitilde100 ns). Communication via shared memory (fast but dangerous). Communication between processes requires IPC (pipes, sockets, shared memory with explicit mapping).

\setcodelang{}

\begin{verbatim}
Thread Creation Cost:         Process Creation Cost:
- Stack allocation            - New address space
- Register state              - Copy page tables (or CoW)
- TCB (Thread Control Block)  - Duplicate file descriptors
≈ microseconds                ≈ milliseconds (fork) to tens of ms
\end{verbatim}

\textbf{When to use processes:} Isolation needed (security, stability --- crash doesn\textquotesingle t kill parent), different languages/runtimes. Example: Nginx worker processes.

\textbf{When to use threads:} Shared state needed, performance-critical (same process, shared cache), I/O overlap within one task. Example: Java web server handling requests.

\subsection{Synchronization}\label{synchronization}

\textbf{The core problem:} Multiple threads accessing shared mutable state without coordination leads to incorrect results.

\setcodelang{}

\begin{verbatim}
Thread 1: READ counter (= 0)
Thread 2: READ counter (= 0)        ← both read 0!
Thread 1: counter = 0 + 1 → WRITE 1
Thread 2: counter = 0 + 1 → WRITE 1  ← lost update!
Result: counter = 1, expected: 2
\end{verbatim}

This is a \textbf{race condition} --- outcome depends on thread scheduling order.

\textbf{Critical Section:} Code that accesses shared state. Must be executed by only one thread at a time (mutual exclusion).

Synchronization mechanisms ranked by cost (cheap \(\rightarrow\) expensive):

\setcodelang{}

\begin{verbatim}
Atomic ops (CAS)  →  Spinlock  →  Mutex  →  Semaphore  →  Monitor
   ~1 ns               ~10 ns      ~25 ns      ~50 ns       ~100 ns
\end{verbatim}

\subsection{Mutex (Mutual Exclusion Lock)}\label{mutex-mutual-exclusion-lock}

A mutex is a binary lock: one thread holds it, all others block (sleep) until it\textquotesingle s released.

\setcodelang{Python}

\begin{Shaded}
\begin{Highlighting}[]
\CommentTok{\# Python example}
\ImportTok{import}\NormalTok{ threading}

\NormalTok{counter }\OperatorTok{=} \DecValTok{0}
\NormalTok{lock }\OperatorTok{=}\NormalTok{ threading.Lock()}

\KeywordTok{def}\NormalTok{ increment():}
    \KeywordTok{global}\NormalTok{ counter}
    \ControlFlowTok{with}\NormalTok{ lock:          }\CommentTok{\# acquire on enter, release on exit}
\NormalTok{        counter }\OperatorTok{+=} \DecValTok{1}    \CommentTok{\# critical section — only 1 thread at a time}

\NormalTok{threads }\OperatorTok{=}\NormalTok{ [threading.Thread(target}\OperatorTok{=}\NormalTok{increment) }\ControlFlowTok{for}\NormalTok{ \_ }\KeywordTok{in} \BuiltInTok{range}\NormalTok{(}\DecValTok{1000}\NormalTok{)]}
\ControlFlowTok{for}\NormalTok{ t }\KeywordTok{in}\NormalTok{ threads: t.start()}
\ControlFlowTok{for}\NormalTok{ t }\KeywordTok{in}\NormalTok{ threads: t.join()}
\CommentTok{\# counter is guaranteed to be 1000}
\end{Highlighting}
\end{Shaded}

\setcodelang{C++}

\begin{Shaded}
\begin{Highlighting}[]
\CommentTok{// C++ example}
\BuiltInTok{std::}\NormalTok{mutex mtx}\OperatorTok{;}
\DataTypeTok{int}\NormalTok{ counter }\OperatorTok{=} \DecValTok{0}\OperatorTok{;}

\DataTypeTok{void}\NormalTok{ increment}\OperatorTok{()} \OperatorTok{\{}
    \BuiltInTok{std::}\NormalTok{lock\_guard}\OperatorTok{\textless{}}\BuiltInTok{std::}\NormalTok{mutex}\OperatorTok{\textgreater{}}\NormalTok{ lock}\OperatorTok{(}\NormalTok{mtx}\OperatorTok{);}  \CommentTok{// RAII — auto releases}
\NormalTok{    counter}\OperatorTok{++;}
\OperatorTok{\}}
\end{Highlighting}
\end{Shaded}

\textbf{Key properties:}

\begin{itemize}
\tightlist
\item
  \textbf{Ownership:} Only the thread that acquired can release (unlike semaphore)
\item
  \textbf{Blocking:} Waiting threads sleep (no CPU waste, but context switch cost)
\item
  \textbf{Non-recursive by default:} Trying to re-lock from same thread \(\rightarrow\) deadlock (use \texttt{recursive\_mutex} if needed)
\end{itemize}

\textbf{Interview trap:} What happens if you forget to release a mutex? \(\rightarrow\) All other threads block forever (deadlock). Always use RAII or try/finally.

\subsection{Semaphore}\label{semaphore}

A semaphore is a counter (initialized to N) with two operations:

\begin{itemize}
\tightlist
\item
  \textbf{wait() / P() / acquire():} Decrement counter. If counter \textless{} 0, block.
\item
  \textbf{signal() / V() / release():} Increment counter. Wake one blocked thread if any.
\end{itemize}

\setcodelang{Python}

\begin{Shaded}
\begin{Highlighting}[]
\NormalTok{Binary Semaphore (N}\OperatorTok{=}\DecValTok{1}\NormalTok{):  behaves like mutex BUT no ownership}
\NormalTok{Counting Semaphore (N}\OperatorTok{=}\NormalTok{k): allows k threads to enter simultaneously}
\end{Highlighting}
\end{Shaded}

\textbf{Use cases:}

\begin{itemize}
\tightlist
\item
  \textbf{Controlling resource pool size:} N=5 \(\rightarrow\) at most 5 threads use DB connection pool
\item
  \textbf{Signaling between threads:} Producer signals, consumer waits
\item
  \textbf{Unlike mutex:} Any thread can signal (not just the acquirer) --- useful for producer/consumer
\end{itemize}

\setcodelang{Java}

\begin{Shaded}
\begin{Highlighting}[]
\CommentTok{// Java: Semaphore for connection pool}
\BuiltInTok{Semaphore}\NormalTok{ pool }\OperatorTok{=} \KeywordTok{new} \BuiltInTok{Semaphore}\OperatorTok{(}\DecValTok{10}\OperatorTok{);} \CommentTok{// max 10 concurrent DB connections}

\DataTypeTok{void} \FunctionTok{queryDB}\OperatorTok{()} \OperatorTok{\{}
\NormalTok{    pool}\OperatorTok{.}\FunctionTok{acquire}\OperatorTok{();}     \CommentTok{// blocks if 10 already in use}
    \ControlFlowTok{try} \OperatorTok{\{}
        \CommentTok{// use DB connection}
    \OperatorTok{\}} \ControlFlowTok{finally} \OperatorTok{\{}
\NormalTok{        pool}\OperatorTok{.}\FunctionTok{release}\OperatorTok{();} \CommentTok{// any thread can release}
    \OperatorTok{\}}
\OperatorTok{\}}
\end{Highlighting}
\end{Shaded}

\subsection{Race Conditions}\label{race-conditions}

A race condition occurs when program correctness depends on the relative timing or ordering of concurrent operations.

\textbf{Three conditions for a race:}

\begin{enumerate}
\def\labelenumi{\arabic{enumi}.}
\tightlist
\item
  Shared mutable state
\item
  At least one writer
\item
  No synchronization
\end{enumerate}

\textbf{Types:}

\begin{itemize}
\tightlist
\item
  \textbf{Read-Write race:} One thread writes while another reads
\item
  \textbf{Write-Write race:} Two threads write simultaneously
\item
  \textbf{Check-Then-Act race:} \texttt{if\ (file\ doesn\textquotesingle{}t\ exist)\ create\ file} --- window between check and act
\end{itemize}

\setcodelang{Java}

\begin{Shaded}
\begin{Highlighting}[]
\CommentTok{// Classic TOCTOU (Time of Check to Time of Use) race:}
\ControlFlowTok{if} \OperatorTok{(!}\NormalTok{map}\OperatorTok{.}\FunctionTok{containsKey}\OperatorTok{(}\NormalTok{key}\OperatorTok{))} \OperatorTok{\{}          \CommentTok{// Thread 1 checks: absent}
    \CommentTok{// Thread 2 also checks: absent}
\NormalTok{    map}\OperatorTok{.}\FunctionTok{put}\OperatorTok{(}\NormalTok{key}\OperatorTok{,} \FunctionTok{computeValue}\OperatorTok{(}\NormalTok{key}\OperatorTok{));}  \CommentTok{// Thread 1 inserts}
    \CommentTok{// Thread 2 also inserts! — duplicate work or overwrite}
\OperatorTok{\}}

\CommentTok{// Fix: atomic putIfAbsent or synchronization}
\NormalTok{map}\OperatorTok{.}\FunctionTok{putIfAbsent}\OperatorTok{(}\NormalTok{key}\OperatorTok{,} \FunctionTok{computeValue}\OperatorTok{(}\NormalTok{key}\OperatorTok{));}
\end{Highlighting}
\end{Shaded}

\textbf{Heisenbugs:} Race conditions that disappear when you add logging/debugging (because the act of observing changes timing). Classic in multi-threaded bugs.

\subsection{Deadlock}\label{deadlock}

\visualfigure{../figures/01-concurrency/02-deadlock.pdf}{Deadlock needs all four Coffman conditions at once; the tell-tale sign is a cycle in the wait-for graph. Break any one condition (e.g. lock ordering) to prevent it.}

\textbf{Definition:} Two or more threads permanently block each other, each waiting for a resource held by another.

\setcodelang{}

\begin{verbatim}
Thread 1:  holds Lock A, waiting for Lock B
Thread 2:  holds Lock B, waiting for Lock A
           ↑ circular dependency — both block forever
\end{verbatim}

\setcodelang{Java}

\begin{Shaded}
\begin{Highlighting}[]
\CommentTok{// Classic deadlock:}
\BuiltInTok{Lock}\NormalTok{ lockA }\OperatorTok{=} \KeywordTok{new} \BuiltInTok{ReentrantLock}\OperatorTok{();}
\BuiltInTok{Lock}\NormalTok{ lockB }\OperatorTok{=} \KeywordTok{new} \BuiltInTok{ReentrantLock}\OperatorTok{();}

\CommentTok{// Thread 1:}
\NormalTok{lockA}\OperatorTok{.}\FunctionTok{lock}\OperatorTok{();}
\CommentTok{// context switch here → Thread 2 runs}
\NormalTok{lockB}\OperatorTok{.}\FunctionTok{lock}\OperatorTok{();} \CommentTok{// blocks! Thread 2 holds lockB}

\CommentTok{// Thread 2:}
\NormalTok{lockB}\OperatorTok{.}\FunctionTok{lock}\OperatorTok{();}
\NormalTok{lockA}\OperatorTok{.}\FunctionTok{lock}\OperatorTok{();} \CommentTok{// blocks! Thread 1 holds lockA}
\CommentTok{// DEADLOCK}
\end{Highlighting}
\end{Shaded}

\subsubsection{The 4 Coffman Conditions (ALL must hold for deadlock)}\label{the-4-coffman-conditions-all-must-hold-for-deadlock}

\textbf{MHNC} --- Mutual exclusion, Hold-and-wait, No preemption, Circular wait

\begin{longtable}[]{@{}
  >{\raggedright\arraybackslash}p{(\columnwidth - 4\tabcolsep) * \real{0.1647}}
  >{\raggedright\arraybackslash}p{(\columnwidth - 4\tabcolsep) * \real{0.4442}}
  >{\raggedright\arraybackslash}p{(\columnwidth - 4\tabcolsep) * \real{0.3911}}@{}}
\toprule\noalign{}
\rowcolor{acc!16}
\begin{minipage}[b]{\linewidth}\raggedright
Condition
\end{minipage} & \begin{minipage}[b]{\linewidth}\raggedright
Meaning
\end{minipage} & \begin{minipage}[b]{\linewidth}\raggedright
Prevention Strategy
\end{minipage} \\
\midrule\noalign{}
\endhead
\bottomrule\noalign{}
\endlastfoot
\textbf{M}utual Exclusion & Resources can\textquotesingle t be shared & Use shareable resources where possible \\
\rowcolor{zebra}
\textbf{H}old-and-Wait & Thread holds resource while waiting for another & Acquire all resources at once (atomic) \\
\textbf{N}o Preemption & Resources can\textquotesingle t be forcibly taken & Allow preemption (rollback and retry) \\
\rowcolor{zebra}
\textbf{C}ircular Wait & Circular chain of threads waiting & \textbf{Impose lock ordering} (most practical!) \\
\end{longtable}

\textbf{Breaking deadlock in practice:}

\begin{enumerate}
\def\labelenumi{\arabic{enumi}.}
\tightlist
\item
  \textbf{Lock ordering:} Always acquire locks in the same global order (e.g., always lock A before B)
\item
  \textbf{Lock timeout:} \texttt{tryLock(timeout)} --- give up and retry
\item
  \textbf{Deadlock detection:} Build a wait-for graph; detect cycles (Java thread dumps, database lock monitors)
\end{enumerate}

\setcodelang{Java}

\begin{Shaded}
\begin{Highlighting}[]
\CommentTok{// Lock ordering fix:}
\CommentTok{// Both threads acquire in same order: lockA then lockB}
\NormalTok{lockA}\OperatorTok{.}\FunctionTok{lock}\OperatorTok{();}
\NormalTok{lockB}\OperatorTok{.}\FunctionTok{lock}\OperatorTok{();}
\CommentTok{// ... use both resources}
\NormalTok{lockB}\OperatorTok{.}\FunctionTok{unlock}\OperatorTok{();}
\NormalTok{lockA}\OperatorTok{.}\FunctionTok{unlock}\OperatorTok{();}
\CommentTok{// No deadlock possible — no circular wait}
\end{Highlighting}
\end{Shaded}

\subsection{Livelock \& Starvation}\label{livelock--starvation}

\textbf{Livelock:} Threads are not blocked but keep changing state in response to each other without making progress. Like two people in a hallway stepping left and right to avoid each other.

\setcodelang{}

\begin{verbatim}
Thread 1: detects conflict → backs off → tries again
Thread 2: detects conflict → backs off → tries again
          Both backing off at the same time, forever
\end{verbatim}

\textbf{Fix:} Add randomized backoff (Ethernet CSMA/CD does this).

\textbf{Starvation:} A thread is perpetually denied access to a resource because other threads keep getting priority. Not deadlock (others make progress), but one thread never does.

\textbf{Fix:} Fair queuing --- FIFO ordering for lock acquisition. Java\textquotesingle s \texttt{ReentrantLock(true)} (fair mode) guarantees FIFO.

\subsection{Atomic Operations \& CAS}\label{atomic-operations--cas}

\textbf{Atomic operation:} Indivisible read-modify-write. Guaranteed by hardware to complete without interruption.

\textbf{Compare-And-Swap (CAS):} The building block of lock-free programming.

\setcodelang{Python}

\begin{Shaded}
\begin{Highlighting}[]
\NormalTok{CAS(memory\_location, expected\_value, new\_value):}
    \ControlFlowTok{if} \OperatorTok{*}\NormalTok{memory\_location }\OperatorTok{==}\NormalTok{ expected\_value:}
        \OperatorTok{*}\NormalTok{memory\_location }\OperatorTok{=}\NormalTok{ new\_value}
        \ControlFlowTok{return}\NormalTok{ true  (success)}
    \ControlFlowTok{else}\NormalTok{:}
        \ControlFlowTok{return}\NormalTok{ false (someone }\ControlFlowTok{else}\NormalTok{ changed it)}
    \OperatorTok{//}\NormalTok{ ALL OF THIS IS ONE ATOMIC HARDWARE INSTRUCTION (CMPXCHG on x86)}
\end{Highlighting}
\end{Shaded}

\setcodelang{Java}

\begin{Shaded}
\begin{Highlighting}[]
\CommentTok{// Java AtomicInteger uses CAS internally:}
\BuiltInTok{AtomicInteger}\NormalTok{ counter }\OperatorTok{=} \KeywordTok{new} \BuiltInTok{AtomicInteger}\OperatorTok{(}\DecValTok{0}\OperatorTok{);}

\CommentTok{// Lock{-}free increment:}
\DataTypeTok{int}\NormalTok{ current}\OperatorTok{,}\NormalTok{ next}\OperatorTok{;}
\ControlFlowTok{do} \OperatorTok{\{}
\NormalTok{    current }\OperatorTok{=}\NormalTok{ counter}\OperatorTok{.}\FunctionTok{get}\OperatorTok{();}
\NormalTok{    next }\OperatorTok{=}\NormalTok{ current }\OperatorTok{+} \DecValTok{1}\OperatorTok{;}
\OperatorTok{\}} \ControlFlowTok{while} \OperatorTok{(!}\NormalTok{counter}\OperatorTok{.}\FunctionTok{compareAndSet}\OperatorTok{(}\NormalTok{current}\OperatorTok{,}\NormalTok{ next}\OperatorTok{));}
\CommentTok{// Retry if another thread changed it between get() and compareAndSet()}
\end{Highlighting}
\end{Shaded}

\textbf{ABA Problem:} CAS sees value A, another thread changes A\(\rightarrow\)B\(\rightarrow\)A, CAS thinks nothing changed. Fix: \textbf{AtomicStampedReference} (add version number alongside value).

\textbf{Lock-free vs wait-free:}

\begin{itemize}
\tightlist
\item
  \textbf{Lock-free:} At least one thread makes progress (no global blocking), but individual threads can starve
\item
  \textbf{Wait-free:} Every thread makes progress in bounded steps (stronger guarantee)
\end{itemize}

\subsection{Memory Visibility \& volatile}\label{memory-visibility--volatile}

\textbf{The hidden problem:} On modern CPUs, each core has its own cache. Thread 1 writes \texttt{x\ =\ 1} to its cache; Thread 2 on another core reads \texttt{x} from its cache and sees \texttt{0}. This is a \textbf{memory visibility} problem, not just a race condition.

\setcodelang{}

\begin{verbatim}
Core 0 Cache: x = 1    Core 1 Cache: x = 0    RAM: x = 0 (stale)
              ↑ write not yet flushed           ↑ Thread 2 reads stale value
\end{verbatim}

\textbf{CPU memory reordering:} CPUs and compilers reorder instructions for performance. What you write in code is not necessarily executed in that order.

\setcodelang{Java}

\begin{Shaded}
\begin{Highlighting}[]
\CommentTok{// Can be reordered by compiler/CPU:}
\NormalTok{x }\OperatorTok{=} \DecValTok{1}\OperatorTok{;}
\NormalTok{ready }\OperatorTok{=} \KeywordTok{true}\OperatorTok{;}  \CommentTok{// another thread might see ready=true but x still 0!}
\end{Highlighting}
\end{Shaded}

\textbf{volatile (Java/C++):} Two guarantees:

\begin{enumerate}
\def\labelenumi{\arabic{enumi}.}
\tightlist
\item
  \textbf{Visibility:} Write is immediately flushed to main memory; read always fetches from main memory
\item
  \textbf{Ordering (Java):} Prevents reordering around volatile access (happens-before guarantee)
\end{enumerate}

\setcodelang{Java}

\begin{Shaded}
\begin{Highlighting}[]
\KeywordTok{volatile} \DataTypeTok{boolean}\NormalTok{ ready }\OperatorTok{=} \KeywordTok{false}\OperatorTok{;}
\DataTypeTok{int}\NormalTok{ x }\OperatorTok{=} \DecValTok{0}\OperatorTok{;}

\CommentTok{// Writer thread:}
\NormalTok{x }\OperatorTok{=} \DecValTok{1}\OperatorTok{;}
\NormalTok{ready }\OperatorTok{=} \KeywordTok{true}\OperatorTok{;}  \CommentTok{// volatile write → x=1 guaranteed visible before ready=true}

\CommentTok{// Reader thread:}
\ControlFlowTok{while} \OperatorTok{(!}\NormalTok{ready}\OperatorTok{)} \OperatorTok{\{\}}  \CommentTok{// volatile read}
\BuiltInTok{System}\OperatorTok{.}\FunctionTok{out}\OperatorTok{.}\FunctionTok{println}\OperatorTok{(}\NormalTok{x}\OperatorTok{);}  \CommentTok{// guaranteed to print 1}
\end{Highlighting}
\end{Shaded}

\textbf{volatile is NOT sufficient for compound operations} (read-modify-write). Use \texttt{AtomicInteger} or mutex for those.

\textbf{Java Memory Model (JMM) happens-before rules:}

\begin{itemize}
\tightlist
\item
  Thread start \(\rightarrow\) actions in started thread
\item
  Lock release \(\rightarrow\) subsequent lock acquire
\item
  Volatile write \(\rightarrow\) subsequent volatile read
\item
  Thread join \(\rightarrow\) actions in joined thread
\end{itemize}

\subsection{Spinlock}\label{spinlock}

A spinlock is a lock where the waiting thread \textbf{busy-waits} (spins in a loop) instead of sleeping.

\setcodelang{C++}

\begin{Shaded}
\begin{Highlighting}[]
\CommentTok{// Spinlock implementation:}
\BuiltInTok{std::}\NormalTok{atomic}\OperatorTok{\textless{}}\DataTypeTok{bool}\OperatorTok{\textgreater{}}\NormalTok{ locked}\OperatorTok{\{}\KeywordTok{false}\OperatorTok{\};}

\DataTypeTok{void}\NormalTok{ acquire}\OperatorTok{()} \OperatorTok{\{}
    \ControlFlowTok{while} \OperatorTok{(}\NormalTok{locked}\OperatorTok{.}\NormalTok{exchange}\OperatorTok{(}\KeywordTok{true}\OperatorTok{))} \OperatorTok{\{}  \CommentTok{// CAS loop}
        \CommentTok{// spin — no sleep, no context switch}
    \OperatorTok{\}}
\OperatorTok{\}}

\DataTypeTok{void}\NormalTok{ release}\OperatorTok{()} \OperatorTok{\{}
\NormalTok{    locked}\OperatorTok{.}\NormalTok{store}\OperatorTok{(}\KeywordTok{false}\OperatorTok{);}
\OperatorTok{\}}
\end{Highlighting}
\end{Shaded}

\begin{longtable}[]{@{}
  >{\raggedright\arraybackslash}p{(\columnwidth - 4\tabcolsep) * \real{0.1619}}
  >{\raggedright\arraybackslash}p{(\columnwidth - 4\tabcolsep) * \real{0.3702}}
  >{\raggedright\arraybackslash}p{(\columnwidth - 4\tabcolsep) * \real{0.4679}}@{}}
\toprule\noalign{}
\rowcolor{acc!16}
\begin{minipage}[b]{\linewidth}\raggedright
Feature
\end{minipage} & \begin{minipage}[b]{\linewidth}\raggedright
Mutex
\end{minipage} & \begin{minipage}[b]{\linewidth}\raggedright
Spinlock
\end{minipage} \\
\midrule\noalign{}
\endhead
\bottomrule\noalign{}
\endlastfoot
Waiting & Thread sleeps (OS schedules out) & Thread spins (burns CPU) \\
\rowcolor{zebra}
Context switch & Yes (\textasciitilde1-10 µs overhead) & No \\
Best for & Long hold times, many threads & Very short critical sections, few threads \\
\rowcolor{zebra}
CPU waste & No (blocked thread uses 0 CPU) & Yes (spinning = 100\% CPU on core) \\
Use in kernel & Yes (interrupts can\textquotesingle t sleep) & Yes \\
\end{longtable}

\textbf{Rule of thumb:} Spinlock wins if critical section is shorter than a context switch (\textasciitilde1 µs). Otherwise, mutex wins.

\subsection{Condition Variables}\label{condition-variables}

Condition variables let threads wait for a specific condition to become true, releasing the associated mutex while waiting.

\setcodelang{Python}

\begin{Shaded}
\begin{Highlighting}[]
\NormalTok{Monitor pattern }\OperatorTok{=}\NormalTok{ Mutex }\OperatorTok{+}\NormalTok{ Condition Variable}
\end{Highlighting}
\end{Shaded}

\setcodelang{Java}

\begin{Shaded}
\begin{Highlighting}[]
\CommentTok{// Java: Producer{-}Consumer with condition variables}
\KeywordTok{import} \ImportTok{java}\OperatorTok{.}\ImportTok{util}\OperatorTok{.}\ImportTok{concurrent}\OperatorTok{.}\ImportTok{locks}\OperatorTok{.*;}

\BuiltInTok{Lock}\NormalTok{ lock }\OperatorTok{=} \KeywordTok{new} \BuiltInTok{ReentrantLock}\OperatorTok{();}
\BuiltInTok{Condition}\NormalTok{ notFull  }\OperatorTok{=}\NormalTok{ lock}\OperatorTok{.}\FunctionTok{newCondition}\OperatorTok{();}
\BuiltInTok{Condition}\NormalTok{ notEmpty }\OperatorTok{=}\NormalTok{ lock}\OperatorTok{.}\FunctionTok{newCondition}\OperatorTok{();}
\BuiltInTok{Queue}\OperatorTok{\textless{}}\BuiltInTok{Integer}\OperatorTok{\textgreater{}}\NormalTok{ queue }\OperatorTok{=} \KeywordTok{new} \BuiltInTok{LinkedList}\OperatorTok{\textless{}\textgreater{}();}
\DataTypeTok{int}\NormalTok{ CAPACITY }\OperatorTok{=} \DecValTok{10}\OperatorTok{;}

\DataTypeTok{void} \FunctionTok{produce}\OperatorTok{(}\DataTypeTok{int}\NormalTok{ item}\OperatorTok{)} \KeywordTok{throws} \BuiltInTok{InterruptedException} \OperatorTok{\{}
\NormalTok{    lock}\OperatorTok{.}\FunctionTok{lock}\OperatorTok{();}
    \ControlFlowTok{try} \OperatorTok{\{}
        \ControlFlowTok{while} \OperatorTok{(}\NormalTok{queue}\OperatorTok{.}\FunctionTok{size}\OperatorTok{()} \OperatorTok{==}\NormalTok{ CAPACITY}\OperatorTok{)}
\NormalTok{            notFull}\OperatorTok{.}\FunctionTok{await}\OperatorTok{();}       \CommentTok{// releases lock, waits}
\NormalTok{        queue}\OperatorTok{.}\FunctionTok{add}\OperatorTok{(}\NormalTok{item}\OperatorTok{);}
\NormalTok{        notEmpty}\OperatorTok{.}\FunctionTok{signal}\OperatorTok{();}         \CommentTok{// wake one consumer}
    \OperatorTok{\}} \ControlFlowTok{finally} \OperatorTok{\{}
\NormalTok{        lock}\OperatorTok{.}\FunctionTok{unlock}\OperatorTok{();}
    \OperatorTok{\}}
\OperatorTok{\}}

\DataTypeTok{void} \FunctionTok{consume}\OperatorTok{()} \KeywordTok{throws} \BuiltInTok{InterruptedException} \OperatorTok{\{}
\NormalTok{    lock}\OperatorTok{.}\FunctionTok{lock}\OperatorTok{();}
    \ControlFlowTok{try} \OperatorTok{\{}
        \ControlFlowTok{while} \OperatorTok{(}\NormalTok{queue}\OperatorTok{.}\FunctionTok{isEmpty}\OperatorTok{())}
\NormalTok{            notEmpty}\OperatorTok{.}\FunctionTok{await}\OperatorTok{();}      \CommentTok{// releases lock, waits}
        \DataTypeTok{int}\NormalTok{ item }\OperatorTok{=}\NormalTok{ queue}\OperatorTok{.}\FunctionTok{poll}\OperatorTok{();}
\NormalTok{        notFull}\OperatorTok{.}\FunctionTok{signal}\OperatorTok{();}          \CommentTok{// wake one producer}
        \ControlFlowTok{return}\NormalTok{ item}\OperatorTok{;}
    \OperatorTok{\}} \ControlFlowTok{finally} \OperatorTok{\{}
\NormalTok{        lock}\OperatorTok{.}\FunctionTok{unlock}\OperatorTok{();}
    \OperatorTok{\}}
\OperatorTok{\}}
\end{Highlighting}
\end{Shaded}

\textbf{Critical: Always use \texttt{while} (not \texttt{if}) when checking condition after await.} Spurious wakeups (thread wakes without being signaled --- real OS phenomenon) and multiple waiters can cause the condition to be false again by the time you execute.

\subsection{Producer-Consumer Problem}\label{producer-consumer-problem}

Classic synchronization problem. Bounded buffer shared between producers (add items) and consumers (remove items).

\textbf{Constraints:}

\begin{enumerate}
\def\labelenumi{\arabic{enumi}.}
\tightlist
\item
  Producer must not add to full buffer
\item
  Consumer must not remove from empty buffer
\item
  Only one thread accesses buffer at a time
\end{enumerate}

\setcodelang{Python}

\begin{Shaded}
\begin{Highlighting}[]
\CommentTok{\# Python: Producer{-}Consumer with semaphores}
\ImportTok{import}\NormalTok{ threading, time, random}

\NormalTok{CAPACITY }\OperatorTok{=} \DecValTok{5}
\NormalTok{mutex }\OperatorTok{=}\NormalTok{ threading.Semaphore(}\DecValTok{1}\NormalTok{)   }\CommentTok{\# binary — protects buffer access}
\NormalTok{empty }\OperatorTok{=}\NormalTok{ threading.Semaphore(CAPACITY)  }\CommentTok{\# slots available to produce}
\NormalTok{full  }\OperatorTok{=}\NormalTok{ threading.Semaphore(}\DecValTok{0}\NormalTok{)         }\CommentTok{\# items available to consume}
\BuiltInTok{buffer} \OperatorTok{=}\NormalTok{ []}

\KeywordTok{def}\NormalTok{ producer():}
    \ControlFlowTok{for}\NormalTok{ i }\KeywordTok{in} \BuiltInTok{range}\NormalTok{(}\DecValTok{10}\NormalTok{):}
\NormalTok{        item }\OperatorTok{=}\NormalTok{ random.randint(}\DecValTok{1}\NormalTok{, }\DecValTok{100}\NormalTok{)}
\NormalTok{        empty.acquire()      }\CommentTok{\# wait for empty slot}
\NormalTok{        mutex.acquire()      }\CommentTok{\# enter critical section}
        \BuiltInTok{buffer}\NormalTok{.append(item)}
        \BuiltInTok{print}\NormalTok{(}\SpecialStringTok{f"Produced }\SpecialCharTok{\{}\NormalTok{item}\SpecialCharTok{\}}\SpecialStringTok{, buffer=}\SpecialCharTok{\{}\BuiltInTok{buffer}\SpecialCharTok{\}}\SpecialStringTok{"}\NormalTok{)}
\NormalTok{        mutex.release()      }\CommentTok{\# exit critical section}
\NormalTok{        full.release()       }\CommentTok{\# signal: one more item available}

\KeywordTok{def}\NormalTok{ consumer():}
    \ControlFlowTok{for}\NormalTok{ i }\KeywordTok{in} \BuiltInTok{range}\NormalTok{(}\DecValTok{10}\NormalTok{):}
\NormalTok{        full.acquire()       }\CommentTok{\# wait for item}
\NormalTok{        mutex.acquire()}
\NormalTok{        item }\OperatorTok{=} \BuiltInTok{buffer}\NormalTok{.pop(}\DecValTok{0}\NormalTok{)}
        \BuiltInTok{print}\NormalTok{(}\SpecialStringTok{f"Consumed }\SpecialCharTok{\{}\NormalTok{item}\SpecialCharTok{\}}\SpecialStringTok{, buffer=}\SpecialCharTok{\{}\BuiltInTok{buffer}\SpecialCharTok{\}}\SpecialStringTok{"}\NormalTok{)}
\NormalTok{        mutex.release()}
\NormalTok{        empty.release()      }\CommentTok{\# signal: one more slot available}

\NormalTok{t1 }\OperatorTok{=}\NormalTok{ threading.Thread(target}\OperatorTok{=}\NormalTok{producer)}
\NormalTok{t2 }\OperatorTok{=}\NormalTok{ threading.Thread(target}\OperatorTok{=}\NormalTok{consumer)}
\NormalTok{t1.start()}\OperatorTok{;}\NormalTok{ t2.start()}
\NormalTok{t1.join()}\OperatorTok{;}\NormalTok{ t2.join()}
\end{Highlighting}
\end{Shaded}

\textbf{Key insight:} The ORDER of semaphore operations matters.

\begin{itemize}
\tightlist
\item
  Producer: \texttt{empty.acquire()} THEN \texttt{mutex.acquire()} (wrong order \(\rightarrow\) deadlock with full buffer)
\item
  Consumer: \texttt{full.acquire()} THEN \texttt{mutex.acquire()}
\end{itemize}

\subsection{Reader-Writer Problem}\label{reader-writer-problem}

Multiple readers can read simultaneously (no conflict). Writers need exclusive access.

\textbf{Priority variants:}

\begin{itemize}
\tightlist
\item
  \textbf{Reader priority:} Readers never wait if any reader is active. Writers can starve.
\item
  \textbf{Writer priority:} New readers wait if a writer is waiting. Fairer.
\end{itemize}

\setcodelang{Java}

\begin{Shaded}
\begin{Highlighting}[]
\CommentTok{// Java ReadWriteLock — production{-}quality solution}
\BuiltInTok{ReadWriteLock}\NormalTok{ rwLock }\OperatorTok{=} \KeywordTok{new} \BuiltInTok{ReentrantReadWriteLock}\OperatorTok{();}
\BuiltInTok{Lock}\NormalTok{ readLock  }\OperatorTok{=}\NormalTok{ rwLock}\OperatorTok{.}\FunctionTok{readLock}\OperatorTok{();}
\BuiltInTok{Lock}\NormalTok{ writeLock }\OperatorTok{=}\NormalTok{ rwLock}\OperatorTok{.}\FunctionTok{writeLock}\OperatorTok{();}

\CommentTok{// Multiple readers can hold simultaneously:}
\DataTypeTok{void} \FunctionTok{read}\OperatorTok{()} \OperatorTok{\{}
\NormalTok{    readLock}\OperatorTok{.}\FunctionTok{lock}\OperatorTok{();}
    \ControlFlowTok{try} \OperatorTok{\{} \CommentTok{/* read shared data */} \OperatorTok{\}}
    \ControlFlowTok{finally} \OperatorTok{\{}\NormalTok{ readLock}\OperatorTok{.}\FunctionTok{unlock}\OperatorTok{();} \OperatorTok{\}}
\OperatorTok{\}}

\CommentTok{// Only one writer, and only when no readers:}
\DataTypeTok{void} \FunctionTok{write}\OperatorTok{()} \OperatorTok{\{}
\NormalTok{    writeLock}\OperatorTok{.}\FunctionTok{lock}\OperatorTok{();}
    \ControlFlowTok{try} \OperatorTok{\{} \CommentTok{/* modify shared data */} \OperatorTok{\}}
    \ControlFlowTok{finally} \OperatorTok{\{}\NormalTok{ writeLock}\OperatorTok{.}\FunctionTok{unlock}\OperatorTok{();} \OperatorTok{\}}
\OperatorTok{\}}
\end{Highlighting}
\end{Shaded}

\textbf{State machine:}

\setcodelang{}

\begin{verbatim}
State: UNLOCKED
  → readLock()  → READING (multiple readers OK)
  → writeLock() → WRITING (exclusive)

State: READING (N readers)
  → readLock()   → READING (N+1 readers)
  → writeLock()  → BLOCKS until all readers done
  → readUnlock() → READING (N-1), or UNLOCKED if N=1

State: WRITING
  → readLock()  → BLOCKS
  → writeLock() → BLOCKS
  → writeUnlock() → UNLOCKED
\end{verbatim}

\textbf{Interview insight:} ReadWriteLock only wins when reads are dominant (\textgreater90\%). With heavy writes, the overhead of tracking readers/writers makes it slower than a plain mutex.

\subsection{Thread Pools}\label{thread-pools}

Creating threads is expensive. Thread pools maintain a set of pre-created worker threads that pick up tasks from a queue.

\setcodelang{}

\begin{verbatim}
                    ┌─────────────────────────┐
Submitted tasks →   │    Task Queue (bounded) │
                    └────────────┬────────────┘
                                 │
              ┌──────────────────┼──────────────────┐
              ↓                  ↓                  ↓
         [Worker T1]        [Worker T2]        [Worker T3]
         (idle/busy)        (idle/busy)        (idle/busy)
\end{verbatim}

\textbf{ThreadPoolExecutor parameters (Java):}

\setcodelang{Java}

\begin{Shaded}
\begin{Highlighting}[]
\BuiltInTok{ThreadPoolExecutor}\NormalTok{ pool }\OperatorTok{=} \KeywordTok{new} \BuiltInTok{ThreadPoolExecutor}\OperatorTok{(}
    \DecValTok{4}\OperatorTok{,}           \CommentTok{// corePoolSize: always{-}on threads}
    \DecValTok{8}\OperatorTok{,}           \CommentTok{// maximumPoolSize: max threads when queue full}
    \DecValTok{60L}\OperatorTok{,}         \CommentTok{// keepAliveTime: idle non{-}core threads die after this}
    \BuiltInTok{TimeUnit}\OperatorTok{.}\FunctionTok{SECONDS}\OperatorTok{,}
    \KeywordTok{new} \BuiltInTok{ArrayBlockingQueue}\OperatorTok{\textless{}\textgreater{}(}\DecValTok{100}\OperatorTok{),}  \CommentTok{// bounded task queue}
    \KeywordTok{new} \BuiltInTok{ThreadPoolExecutor}\OperatorTok{.}\FunctionTok{CallerRunsPolicy}\OperatorTok{()}  \CommentTok{// rejection policy}
\OperatorTok{);}
\end{Highlighting}
\end{Shaded}

\textbf{Rejection policies when queue full + max threads reached:}

\begin{itemize}
\tightlist
\item
  \texttt{AbortPolicy} (default): throw \texttt{RejectedExecutionException}
\item
  \texttt{CallerRunsPolicy}: caller thread runs the task (backpressure)
\item
  \texttt{DiscardPolicy}: silently drop the task
\item
  \texttt{DiscardOldestPolicy}: drop oldest queued task, retry
\end{itemize}

\textbf{Pool sizing formula:}

\setcodelang{}

\begin{verbatim}
CPU-bound tasks:   threads ≈ CPU_cores + 1 (extra for occasional blocking)
I/O-bound tasks:   threads ≈ CPU_cores × (1 + wait_time/compute_time)
Example: 8 cores, 90% I/O wait → 8 × (1 + 9) = 80 threads
\end{verbatim}

\subsection{False Sharing}\label{false-sharing}

Two threads modify different variables that happen to sit in the \textbf{same cache line} (64 bytes on x86). Each core\textquotesingle s write invalidates the other core\textquotesingle s cache line, causing constant cache misses --- as slow as sharing a variable, but no actual sharing.

\setcodelang{}

\begin{verbatim}
Cache line (64 bytes):
[counter_A (8 bytes)][counter_B (8 bytes)][... padding ...]
     ↑ Thread 1            ↑ Thread 2
     writes here           writes here
     → invalidates Thread 2's copy of ENTIRE cache line
     → Thread 2 must reload from memory
     → thrashing even though they access different variables!
\end{verbatim}

\setcodelang{Java}

\begin{Shaded}
\begin{Highlighting}[]
\CommentTok{// Fix: padding to separate variables into different cache lines}
\KeywordTok{class}\NormalTok{ PaddedCounter }\OperatorTok{\{}
    \DataTypeTok{long}\NormalTok{ p1}\OperatorTok{,}\NormalTok{ p2}\OperatorTok{,}\NormalTok{ p3}\OperatorTok{,}\NormalTok{ p4}\OperatorTok{,}\NormalTok{ p5}\OperatorTok{,}\NormalTok{ p6}\OperatorTok{,}\NormalTok{ p7}\OperatorTok{;} \CommentTok{// 56 bytes padding}
    \KeywordTok{volatile} \DataTypeTok{long}\NormalTok{ value}\OperatorTok{;}               \CommentTok{// 8 bytes — own cache line}
    \DataTypeTok{long}\NormalTok{ q1}\OperatorTok{,}\NormalTok{ q2}\OperatorTok{,}\NormalTok{ q3}\OperatorTok{,}\NormalTok{ q4}\OperatorTok{,}\NormalTok{ q5}\OperatorTok{,}\NormalTok{ q6}\OperatorTok{,}\NormalTok{ q7}\OperatorTok{;} \CommentTok{// 56 bytes padding}
\OperatorTok{\}}
\CommentTok{// Or use @Contended annotation in Java 8+ (JEP 142)}
\end{Highlighting}
\end{Shaded}

\textbf{LMAX Disruptor} --- high-performance ring buffer --- was designed specifically to eliminate false sharing.

\subsection{Python GIL (Global Interpreter Lock)}\label{python-gil-global-interpreter-lock}

CPython (the standard Python interpreter) has a global mutex --- the GIL --- that only allows one thread to execute Python bytecode at a time.

\setcodelang{}

\begin{verbatim}
Thread 1: [Python bytecode]─┐
Thread 2:                   └─ GIL allows only one at a time
Thread 3: [C extension]────── (bypasses GIL — can run in parallel!)
\end{verbatim}

\textbf{Implications:}

\begin{itemize}
\tightlist
\item
  Python threads don\textquotesingle t give CPU parallelism for CPU-bound work
\item
  Threads still help for I/O-bound work (GIL released during I/O)
\item
  For CPU parallelism: use \texttt{multiprocessing} (separate processes, each with own GIL) or C extensions (NumPy, TensorFlow bypass the GIL)
\end{itemize}

\setcodelang{Python}

\begin{Shaded}
\begin{Highlighting}[]
\CommentTok{\# CPU{-}bound: threads DON\textquotesingle{}T help (GIL bottleneck)}
\ImportTok{import}\NormalTok{ threading}
\CommentTok{\# 4 threads still run sequentially for pure Python computation}

\CommentTok{\# CPU{-}bound: use multiprocessing}
\ImportTok{from}\NormalTok{ multiprocessing }\ImportTok{import}\NormalTok{ Pool}
\ControlFlowTok{with}\NormalTok{ Pool(}\DecValTok{4}\NormalTok{) }\ImportTok{as}\NormalTok{ p:}
\NormalTok{    results }\OperatorTok{=}\NormalTok{ p.}\BuiltInTok{map}\NormalTok{(cpu\_heavy\_fn, data)  }\CommentTok{\# true parallelism}

\CommentTok{\# I/O{-}bound: threads work fine}
\ImportTok{import}\NormalTok{ threading}
\CommentTok{\# threads waiting on network/disk release GIL → true concurrency}
\end{Highlighting}
\end{Shaded}

\textbf{Python 3.13+:} "No-GIL" build (PEP 703) is in experimental phase, using fine-grained per-object locking instead.

\subsection{Async vs Threads}\label{async-vs-threads}

Both deal with concurrency, but with different execution models.

\setcodelang{Python}

\begin{Shaded}
\begin{Highlighting}[]
\NormalTok{Threads:                              Async (coroutines):}
\NormalTok{OS schedules preemptively             Cooperative — yields explicitly}
\NormalTok{Each thread has own stack (MB)        One stack shared, tiny frames (KB)}
\NormalTok{Context switch: OS interrupt          Context switch: }\ControlFlowTok{await}\NormalTok{ keyword}
\NormalTok{Good }\ControlFlowTok{for}\NormalTok{: CPU}\OperatorTok{{-}}\NormalTok{bound, blocking C libs  Good }\ControlFlowTok{for}\NormalTok{: I}\OperatorTok{/}\NormalTok{O}\OperatorTok{{-}}\NormalTok{bound, many connections}

\NormalTok{Node.js}\OperatorTok{/}\NormalTok{Python asyncio }\OperatorTok{=} \DecValTok{1}\NormalTok{ thread, event loop, callbacks}\OperatorTok{/}\NormalTok{coroutines}
\NormalTok{Java Virtual Threads (Project Loom) }\OperatorTok{=} \ControlFlowTok{async}\NormalTok{ benefits }\ControlFlowTok{with}\NormalTok{ thread API}
\end{Highlighting}
\end{Shaded}

\setcodelang{Python}

\begin{Shaded}
\begin{Highlighting}[]
\CommentTok{\# Async Python: 10,000 concurrent HTTP requests, 1 thread}
\ImportTok{import}\NormalTok{ asyncio, aiohttp}

\ControlFlowTok{async} \KeywordTok{def}\NormalTok{ fetch(session, url):}
    \ControlFlowTok{async} \ControlFlowTok{with}\NormalTok{ session.get(url) }\ImportTok{as}\NormalTok{ response:}
        \ControlFlowTok{return} \ControlFlowTok{await}\NormalTok{ response.text()}

\ControlFlowTok{async} \KeywordTok{def}\NormalTok{ main():}
    \ControlFlowTok{async} \ControlFlowTok{with}\NormalTok{ aiohttp.ClientSession() }\ImportTok{as}\NormalTok{ session:}
\NormalTok{        tasks }\OperatorTok{=}\NormalTok{ [fetch(session, url) }\ControlFlowTok{for}\NormalTok{ url }\KeywordTok{in}\NormalTok{ urls]}
\NormalTok{        results }\OperatorTok{=} \ControlFlowTok{await}\NormalTok{ asyncio.gather(}\OperatorTok{*}\NormalTok{tasks)}
\end{Highlighting}
\end{Shaded}

\textbf{The fundamental difference:}

\setcodelang{}

\begin{verbatim}
Threads: preemptive multitasking — OS can switch anytime
Async:   cooperative multitasking — switch only at await points
         → no race conditions between await points (within one coroutine)
         → but: long-running sync code blocks the entire event loop
\end{verbatim}

\section{Architecture / Diagrams}\label{architecture--diagrams}

\subsection{Thread Lifecycle}\label{thread-lifecycle}

\setcodelang{}

\begin{verbatim}
                    new Thread()
                         │
                         ▼
                     [NEW/CREATED]
                         │
                    thread.start()
                         │
                         ▼
                    [RUNNABLE]◄────────────────────┐
                         │                         │
              ┌──────────┼──────────┐              │
              │          │          │              │
         lock busy   I/O wait    sleep()        wakeup/
              │          │          │           timeout
              ▼          ▼          ▼              │
          [BLOCKED]  [WAITING]  [TIMED_WAITING]───┘
                                                   
                    run() returns / exception
                         │
                         ▼
                    [TERMINATED]
\end{verbatim}

\subsection{Mutex vs Semaphore vs Monitor}\label{mutex-vs-semaphore-vs-monitor}

\setcodelang{}

\begin{verbatim}
┌──────────────┬──────────────────────┬────────────────────────────┐
│              │ Mutex                │ Semaphore                  │
├──────────────┼──────────────────────┼────────────────────────────┤
│ Value        │ binary (0 or 1)      │ integer (0 to N)           │
│ Ownership    │ yes (locked thread   │ no (any thread can signal) │
│              │ must unlock)         │                            │
│ Use case     │ mutual exclusion     │ signaling, resource count  │
│ Recursive?   │ no (by default)      │ yes (counter-based)        │
│ Priority     │ can do inheritance   │ typically not              │
└──────────────┴──────────────────────┴────────────────────────────┘

Monitor = Object/Class with:
  - Internal mutex (one thread at a time)
  - Condition variables (wait/notify)
  Java's synchronized methods/blocks = Monitor
\end{verbatim}

\subsection{Deadlock Detection: Wait-For Graph}\label{deadlock-detection-wait-for-graph}

\setcodelang{}

\begin{verbatim}
T1 ──waits──► T2
▲              │
│              waits
└──waits── T3 ◄┘

Cycle T1→T2→T3→T1 = DEADLOCK
\end{verbatim}

\subsection{Lock Ordering to Prevent Deadlock}\label{lock-ordering-to-prevent-deadlock}

\setcodelang{}

\begin{verbatim}
Global lock order: L1 < L2 < L3 < L4

Thread A needs L2, L4:   acquire L2 → acquire L4 → release L4 → release L2
Thread B needs L1, L4:   acquire L1 → acquire L4 → release L4 → release L1
Thread C needs L2, L3:   acquire L2 → acquire L3 → release L3 → release L2

No circular dependency possible — all follow same ordering
\end{verbatim}

\subsection{CAS Operation Internals}\label{cas-operation-internals}

\setcodelang{Python}

\begin{Shaded}
\begin{Highlighting}[]
\NormalTok{CMPXCHG instruction (x86):}
\FloatTok{1.}\NormalTok{ Read value }\ImportTok{from}\NormalTok{ memory address}
\FloatTok{2.}\NormalTok{ Compare }\ControlFlowTok{with}\NormalTok{ expected}
\FloatTok{3.}\NormalTok{ If equal: write new value, }\ControlFlowTok{return}\NormalTok{ true}
\FloatTok{4.}\NormalTok{ If }\KeywordTok{not}\NormalTok{ equal: }\ControlFlowTok{return}\NormalTok{ false (read actual value)}
\NormalTok{Steps }\DecValTok{1}\OperatorTok{{-}}\DecValTok{4}\NormalTok{ are ATOMIC — cannot be interrupted by another CPU}
\end{Highlighting}
\end{Shaded}

\subsection{Memory Model: Visibility Issue}\label{memory-model-visibility-issue}

\setcodelang{}

\begin{verbatim}
Multi-core CPU:

Core 0                    Core 1
┌──────┐                 ┌──────┐
│ L1 $ │ x=1 (dirty)     │ L1 $ │ x=0 (stale)
└──┬───┘                 └──┬───┘
   │                        │
┌──▼────────────────────────▼──┐
│        L3 Cache / RAM         │  x=0 (not yet flushed)
└───────────────────────────────┘

volatile/memory fence forces:
- Write: flush to RAM immediately
- Read: invalidate local cache, read from RAM
\end{verbatim}

\subsection{Thread Pool Flow}\label{thread-pool-flow}

\setcodelang{}

\begin{verbatim}
Client threads
     │
     ▼
[submit(Runnable)] ──► [BlockingQueue] ──► [Worker Thread Pool]
                         capacity=N         core=4, max=8
                              │
                         queue full?
                              │
                    ┌─────────┴──────────┐
                    │                    │
             spawn extra worker    queue still full?
             (up to maxPoolSize)         │
                                  [Rejection Policy]
\end{verbatim}

\section{Real-World Examples}\label{real-world-examples}

\subsection{Google: Bigtable Tablet Server}\label{google-bigtable-tablet-server}

Bigtable tablet servers use read-write locks to protect tablet metadata. Multiple reads (serving queries) proceed concurrently; mutations (compactions, splits) acquire write lock exclusively. This allows high read throughput while ensuring consistency during structural changes.

\subsection{Meta: Folly\textquotesingle s Concurrency Primitives}\label{meta-follys-concurrency-primitives}

Meta\textquotesingle s open-source Folly library (\texttt{folly::Synchronized\textless{}T\textgreater{}}) wraps any data structure with a mutex automatically. Their \texttt{SharedMutex} is a high-performance reader-writer lock used in Thrift servers where millions of RPC calls read configuration but writes are rare.

\subsection{Amazon DynamoDB: Lock-Free Data Structures}\label{amazon-dynamodb-lock-free-data-structures}

DynamoDB uses lock-free concurrent skip lists for its in-memory data structures. CAS operations handle concurrent inserts/deletes without traditional locking, achieving high throughput at low latency for partition-local operations.

\subsection{Uber: Atomic Counters for Surge Pricing}\label{uber-atomic-counters-for-surge-pricing}

Uber\textquotesingle s surge pricing uses \texttt{AtomicLong} counters to track active trips per geohash cell. When a driver completes a trip, \texttt{decrementAndGet()} happens atomically across multiple dispatch servers using CAS --- no locks needed for this critical counting operation.

\subsection{Linux Kernel: RCU (Read-Copy-Update)}\label{linux-kernel-rcu-read-copy-update}

RCU is an extreme read-write lock optimization: readers never lock at all (zero overhead). Writers make a copy, update it, then atomically swap the pointer. Old copy is freed when no readers reference it. Used in Linux networking, file system paths --- thousands of reads per microsecond.

\subsection{Java\textquotesingle s ConcurrentHashMap}\label{javas-concurrenthashmap}

Pre-Java 8: one lock per segment (16 segments = 16-way parallelism). Java 8+: CAS for \texttt{put} when bucket is empty; per-bucket \texttt{synchronized} when needed. Readers never block. This gives near-linear scalability with thread count for most workloads.

\section{Real-Life Analogies}\label{real-life-analogies}

\emph{One kitchen, one brigade --- every concept is a cook, a tool, or a station in the same busy restaurant.}

\begin{longtable}[]{@{}
  >{\raggedright\arraybackslash}p{(\columnwidth - 2\tabcolsep) * \real{0.1962}}
  >{\raggedright\arraybackslash}p{(\columnwidth - 2\tabcolsep) * \real{0.8038}}@{}}
\toprule\noalign{}
\rowcolor{acc!16}
\begin{minipage}[b]{\linewidth}\raggedright
Concept
\end{minipage} & \begin{minipage}[b]{\linewidth}\raggedright
Analogy --- \emph{the restaurant kitchen}
\end{minipage} \\
\midrule\noalign{}
\endhead
\bottomrule\noalign{}
\endlastfoot
\textbf{Thread} & A line cook at a station --- shares the kitchen (process) but works their own tickets and prep (stack) \\
\rowcolor{zebra}
\textbf{Process} & A separate restaurant down the street --- its own kitchen and pantry; the two only talk by phone (IPC) \\
\textbf{Mutex} & The single razor-sharp chef\textquotesingle s knife --- only one cook can hold it; the rest wait their turn to chop \\
\rowcolor{zebra}
\textbf{Semaphore} & The stove\textquotesingle s 4 burners --- at most 4 dishes cook at once; finish one and the next cook claims the burner \\
\textbf{Deadlock} & Cook A holds the only pan and needs the oven; Cook B holds the oven and needs the pan --- both freeze \\
\rowcolor{zebra}
\textbf{Livelock} & Two cooks meet in the narrow doorway and keep stepping the same way to yield --- polite, but nobody passes \\
\textbf{Starvation} & The new cook who\textquotesingle s always skipped for the busy grill --- everyone else keeps getting picked first \\
\rowcolor{zebra}
\textbf{Race Condition} & Two cooks grab the last egg for different orders --- one ticket goes out wrong \\
\textbf{Spinlock} & A cook hovering at the oven, opening it every few seconds to check if it\textquotesingle s free instead of prepping \\
\rowcolor{zebra}
\textbf{Condition Variable} & The expo bell --- cooks rest until "order up!" rings, then wake and plate \\
\textbf{volatile} & The shared specials whiteboard --- everyone reads the same board instantly, not private notepads \\
\rowcolor{zebra}
\textbf{CAS} & "Is table 5\textquotesingle s ticket still unclaimed?" --- glance at the rail before grabbing it; if taken, try the next \\
\textbf{Thread Pool} & A fixed brigade pulling tickets off the rail --- idle cooks wait, finished cooks grab the next ticket \\
\rowcolor{zebra}
\textbf{False Sharing} & Two cooks elbowing over one cutting board for different vegetables --- both slow down without truly sharing \\
\textbf{ABA Problem} & A burner looked free; a cook used and cleared it while you blinked --- you assume nothing changed, get burned \\
\rowcolor{zebra}
\textbf{Producer-Consumer} & Chefs plate onto the pass shelf (limited space); waiters carry out --- chefs pause when full, waiters wait when empty \\
\textbf{Reader-Writer} & The recipe board --- any number of cooks read at once, but the head chef locks it alone to rewrite a recipe \\
\rowcolor{zebra}
\textbf{GIL} & A tiny kitchen with a single stove --- hire ten cooks, but only one can actually cook at any moment \\
\textbf{Async/Await} & One waiter working ten tables --- takes the next order while the kitchen cooks the last, never idle \\
\end{longtable}

\section{Memory Tricks / Mnemonics}\label{memory-tricks--mnemonics}

\subsection{\texorpdfstring{Deadlock 4 Conditions: \textbf{"MHNC"} (Make Him Never Circulate)}{Deadlock 4 Conditions: "MHNC" (Make Him Never Circulate)}}\label{deadlock-4-conditions-mhnc-make-him-never-circulate}

\begin{itemize}
\tightlist
\item
  \textbf{M}utual Exclusion
\item
  \textbf{H}old and Wait
\item
  \textbf{N}o Preemption
\item
  \textbf{C}ircular Wait
\end{itemize}

\textbf{Break any ONE to prevent deadlock:}

\begin{itemize}
\tightlist
\item
  Break M: Use shareable resources (read-only data)
\item
  Break H: Acquire all locks atomically at once
\item
  Break N: Allow lock preemption (databases do this with rollback)
\item
  Break C: \textbf{LOCK ORDERING} --- most practical solution in software
\end{itemize}

\subsection{\texorpdfstring{Semaphore P/V: \textbf{"P = Post office, V = Vacate"}}{Semaphore P/V: "P = Post office, V = Vacate"}}\label{semaphore-pv-p--post-office-v--vacate}

\begin{itemize}
\tightlist
\item
  \textbf{P} (proberen/test) = acquire, wait, decrement \(\rightarrow\) "waiting in POST office queue"
\item
  \textbf{V} (verhogen/increment) = release, signal \(\rightarrow\) "VACATE the window, next!"
\end{itemize}

\subsection{\texorpdfstring{Thread States: \textbf{"NEW RABBIT WAITS BUSILY THEN DIES"}}{Thread States: "NEW RABBIT WAITS BUSILY THEN DIES"}}\label{thread-states-new-rabbit-waits-busily-then-dies}

\begin{itemize}
\tightlist
\item
  NEW \(\rightarrow\) RUNNABLE \(\rightarrow\) WAITING/TIMED\_WAITING/BLOCKED \(\rightarrow\) TERMINATED
\end{itemize}

\subsection{\texorpdfstring{Happens-Before Rules (Java): \textbf{"SLVJ"} --- Start, Lock, Volatile, Join}{Happens-Before Rules (Java): "SLVJ" --- Start, Lock, Volatile, Join}}\label{happens-before-rules-java-slvj--start-lock-volatile-join}

\begin{itemize}
\tightlist
\item
  \textbf{S}tart: thread.start() happens-before thread\textquotesingle s actions
\item
  \textbf{L}ock: unlock happens-before subsequent lock
\item
  \textbf{V}olatile: write happens-before subsequent read
\item
  \textbf{J}oin: thread actions happen-before join() returns
\end{itemize}

\subsection{\texorpdfstring{CAS vs Lock: \textbf{"CAS for counters, Lock for compounds"}}{CAS vs Lock: "CAS for counters, Lock for compounds"}}\label{cas-vs-lock-cas-for-counters-lock-for-compounds}

\begin{itemize}
\tightlist
\item
  Single variable update \(\rightarrow\) CAS/Atomic
\item
  Multiple variables must change together \(\rightarrow\) Lock
\end{itemize}

\subsection{\texorpdfstring{Pool Sizing: \textbf{"I/O = Cores \(\times\) 10, CPU = Cores + 1"}}{Pool Sizing: "I/O = Cores \textbackslash times 10, CPU = Cores + 1"}}\label{pool-sizing-io--cores-times-10-cpu--cores--1}

\begin{itemize}
\tightlist
\item
  Pure CPU work: \texttt{cores\ +\ 1} threads
\item
  Pure I/O work: \texttt{cores\ ×\ 10} threads (rough rule)
\item
  Mixed: measure and tune with load testing
\end{itemize}

\subsection{\texorpdfstring{False Sharing Fix: \textbf{"Pad to 64"}}{False Sharing Fix: "Pad to 64"}}\label{false-sharing-fix-pad-to-64}

\begin{itemize}
\tightlist
\item
  Cache line = 64 bytes on x86
\item
  Pad hot variables so they don\textquotesingle t share a cache line
\end{itemize}

\subsection{volatile vs synchronized:}\label{volatile-vs-synchronized}

\begin{itemize}
\tightlist
\item
  \textbf{volatile:} single-variable visibility + ordering (no atomicity for compound ops)
\item
  \textbf{synchronized:} visibility + ordering + atomicity for critical section
\end{itemize}

\subsection{\texorpdfstring{Memory model mental model: \textbf{"Write \(\rightarrow\) Fence \(\rightarrow\) Read"}}{Memory model mental model: "Write \textbackslash rightarrow Fence \textbackslash rightarrow Read"}}\label{memory-model-mental-model-write-rightarrow-fence-rightarrow-read}

\begin{itemize}
\tightlist
\item
  volatile write = release fence (all prior writes visible after)
\item
  volatile read = acquire fence (all subsequent reads see previous writes)
\end{itemize}

\section{Common Interview Questions}\label{common-interview-questions}

\subsection{Q1: What is a race condition? How do you prevent it?}\label{q1-what-is-a-race-condition-how-do-you-prevent-it}

\textbf{Model answer:} A race condition occurs when program output depends on relative thread scheduling. Prevents: identify shared mutable state, protect with mutex/synchronized, or use lock-free atomics, or eliminate sharing (immutability, thread-local storage).

\textbf{Follow-ups:}

\begin{itemize}
\tightlist
\item
  "Give me an example of a race condition in a real system" \(\rightarrow\) Double-checked locking without volatile in pre-Java 5 singleton
\item
  "Can you have a race condition with a single CPU?" \(\rightarrow\) Yes, due to context switches between instructions
\end{itemize}

\subsection{Q2: Explain the four conditions for deadlock and how to prevent each.}\label{q2-explain-the-four-conditions-for-deadlock-and-how-to-prevent-each}

\textbf{Model answer:} MHNC --- Mutual Exclusion, Hold-and-Wait, No Preemption, Circular Wait. Prevention: lock ordering (breaks circular wait), acquire all locks at once (breaks hold-and-wait), tryLock with timeout (breaks no-preemption).

\textbf{Follow-ups:}

\begin{itemize}
\tightlist
\item
  "How does a database detect deadlock?" \(\rightarrow\) Wait-for graph cycle detection; victim selection by transaction age/cost
\item
  "What\textquotesingle s the difference between deadlock prevention and avoidance?" \(\rightarrow\) Prevention breaks a condition; avoidance (Banker\textquotesingle s algorithm) dynamically refuses unsafe states
\end{itemize}

\subsection{Q3: Implement producer-consumer with two semaphores.}\label{q3-implement-producer-consumer-with-two-semaphores}

\textbf{Model answer:} (See code in §Producer-Consumer section above)

\textbf{Key points to mention:}

\begin{itemize}
\tightlist
\item
  \texttt{empty} semaphore initialized to buffer capacity
\item
  \texttt{full} semaphore initialized to 0
\item
  Acquire resource semaphore BEFORE mutex to avoid deadlock
\item
  Why while instead of if for condition check
\end{itemize}

\textbf{Follow-ups:}

\begin{itemize}
\tightlist
\item
  "What if we have 3 producers and 5 consumers?" \(\rightarrow\) Same code, semaphores handle N:M automatically
\item
  "How would you implement this without semaphores?" \(\rightarrow\) Mutex + condition variables (as in Java example)
\end{itemize}

\subsection{Q4: What is the difference between a mutex and a semaphore?}\label{q4-what-is-the-difference-between-a-mutex-and-a-semaphore}

\textbf{Model answer:}

\begin{itemize}
\tightlist
\item
  Mutex: binary, has ownership (acquirer must release), used for mutual exclusion
\item
  Semaphore: counter 0-N, no ownership (any thread can signal), used for resource counting and signaling
\end{itemize}

\textbf{Follow-up:} "Can you implement a mutex with a semaphore?" \(\rightarrow\) Yes, binary semaphore with convention that only acquirer calls signal. But you lose deadlock detection and priority inheritance benefits.

\subsection{Q5: Explain Compare-And-Swap and the ABA problem.}\label{q5-explain-compare-and-swap-and-the-aba-problem}

\textbf{Model answer:} CAS atomically: reads value, compares to expected, writes new value if match, returns success/fail. Single hardware instruction (CMPXCHG). ABA: value changed A\(\rightarrow\)B\(\rightarrow\)A between your read and CAS --- CAS succeeds thinking nothing changed. Fix: add version counter (\texttt{AtomicStampedReference}).

\textbf{Follow-up:} "Where would ABA actually cause a bug?" \(\rightarrow\) Lock-free linked list: node removed, memory freed, new node allocated at same address. CAS sees same pointer, thinks old node still there \(\rightarrow\) use-after-free or corruption.

\subsection{Q6: What is false sharing and how do you fix it?}\label{q6-what-is-false-sharing-and-how-do-you-fix-it}

\textbf{Model answer:} Two threads write to different variables sharing a 64-byte cache line. Each write invalidates the other core\textquotesingle s cache line, causing cache misses even though the data is logically independent. Fix: pad struct so hot variables occupy their own cache line.

\textbf{Follow-up:} "How do you detect false sharing?" \(\rightarrow\) \texttt{perf\ stat\ -e\ cache-misses} on Linux, Intel VTune "false sharing" analysis. Look for high L1/L2 cache miss rates with no logical sharing.

\subsection{Q7: What is the Python GIL and when does it matter?}\label{q7-what-is-the-python-gil-and-when-does-it-matter}

\textbf{Model answer:} Global Interpreter Lock --- CPython mutex preventing multiple native threads from executing Python bytecode simultaneously. Matters for CPU-bound parallelism (threads don\textquotesingle t help). Doesn\textquotesingle t matter for I/O-bound concurrency (GIL released during I/O, threads work fine). Fix for CPU-bound: \texttt{multiprocessing}, C extensions, or PyPy.

\textbf{Follow-up:} "Is the GIL being removed?" \(\rightarrow\) Python 3.13 experimental no-GIL build (PEP 703). Not default yet. Uses per-object reference counting with atomic ops instead.

\subsection{Q8: Explain volatile in Java. When is it sufficient?}\label{q8-explain-volatile-in-java-when-is-it-sufficient}

\textbf{Model answer:} Volatile ensures visibility (writes immediately visible to other threads via memory barrier) and ordering (prevents reordering around the volatile access). Sufficient for: flag variables, single-variable state published once. NOT sufficient for: \texttt{counter++} (read-modify-write is not atomic). Use \texttt{AtomicInteger} for that.

\textbf{Follow-up:} "What\textquotesingle s the difference between volatile and synchronized?" \(\rightarrow\) synchronized adds atomicity for compound actions + mutual exclusion. volatile is lighter but no mutual exclusion.

\subsection{Q9: How does Java\textquotesingle s ReadWriteLock work, and when would you not use it?}\label{q9-how-does-javas-readwritelock-work-and-when-would-you-not-use-it}

\textbf{Model answer:} Allows multiple concurrent readers OR one exclusive writer. Writer blocks when readers active; readers block when writer active. Avoids when: write-heavy workloads (overhead \textgreater{} benefit), very short critical sections (overhead of tracking readers), when you need upgrade from read to write lock (causes deadlock --- not supported).

\textbf{Follow-up:} "What\textquotesingle s StampedLock?" \(\rightarrow\) Java 8 addition. Adds optimistic reading: try a read without lock, check stamp at end; if stamp changed, fall back to real read lock. Even better read throughput when conflicts rare.

\subsection{Q10: Design a thread-safe singleton in Java.}\label{q10-design-a-thread-safe-singleton-in-java}

\textbf{Model answer (multiple approaches):}

\setcodelang{Java}

\begin{Shaded}
\begin{Highlighting}[]
\CommentTok{// Approach 1: Eager initialization (simplest, safe)}
\KeywordTok{public} \KeywordTok{class}\NormalTok{ Singleton }\OperatorTok{\{}
    \KeywordTok{private} \DataTypeTok{static} \DataTypeTok{final}\NormalTok{ Singleton INSTANCE }\OperatorTok{=} \KeywordTok{new} \FunctionTok{Singleton}\OperatorTok{();}
    \KeywordTok{public} \DataTypeTok{static}\NormalTok{ Singleton }\FunctionTok{getInstance}\OperatorTok{()} \OperatorTok{\{} \ControlFlowTok{return}\NormalTok{ INSTANCE}\OperatorTok{;} \OperatorTok{\}}
\OperatorTok{\}}

\CommentTok{// Approach 2: Double{-}checked locking (lazy, safe in Java 5+)}
\KeywordTok{public} \KeywordTok{class}\NormalTok{ Singleton }\OperatorTok{\{}
    \KeywordTok{private} \DataTypeTok{static} \KeywordTok{volatile}\NormalTok{ Singleton instance}\OperatorTok{;}  \CommentTok{// volatile REQUIRED}
    \KeywordTok{public} \DataTypeTok{static}\NormalTok{ Singleton }\FunctionTok{getInstance}\OperatorTok{()} \OperatorTok{\{}
        \ControlFlowTok{if} \OperatorTok{(}\NormalTok{instance }\OperatorTok{==} \KeywordTok{null}\OperatorTok{)} \OperatorTok{\{}                  \CommentTok{// first check (no lock)}
            \KeywordTok{synchronized} \OperatorTok{(}\NormalTok{Singleton}\OperatorTok{.}\FunctionTok{class}\OperatorTok{)} \OperatorTok{\{}
                \ControlFlowTok{if} \OperatorTok{(}\NormalTok{instance }\OperatorTok{==} \KeywordTok{null}\OperatorTok{)} \OperatorTok{\{}          \CommentTok{// second check (with lock)}
\NormalTok{                    instance }\OperatorTok{=} \KeywordTok{new} \FunctionTok{Singleton}\OperatorTok{();}
                \OperatorTok{\}}
            \OperatorTok{\}}
        \OperatorTok{\}}
        \ControlFlowTok{return}\NormalTok{ instance}\OperatorTok{;}
    \OperatorTok{\}}
\OperatorTok{\}}

\CommentTok{// Approach 3: Initialization{-}on{-}demand holder (best — lazy + safe)}
\KeywordTok{public} \KeywordTok{class}\NormalTok{ Singleton }\OperatorTok{\{}
    \KeywordTok{private} \DataTypeTok{static} \KeywordTok{class} \BuiltInTok{Holder} \OperatorTok{\{}
        \DataTypeTok{static} \DataTypeTok{final}\NormalTok{ Singleton INSTANCE }\OperatorTok{=} \KeywordTok{new} \FunctionTok{Singleton}\OperatorTok{();}
    \OperatorTok{\}}
    \KeywordTok{public} \DataTypeTok{static}\NormalTok{ Singleton }\FunctionTok{getInstance}\OperatorTok{()} \OperatorTok{\{} \ControlFlowTok{return} \BuiltInTok{Holder}\OperatorTok{.}\FunctionTok{INSTANCE}\OperatorTok{;} \OperatorTok{\}}
\OperatorTok{\}}
\end{Highlighting}
\end{Shaded}

\textbf{Why volatile in DCL?} Without volatile, another thread may see partially constructed object (constructor writes reordered to after \texttt{instance\ =\ new\ Singleton()} write).

\subsection{\texorpdfstring{Q11: What is a condition variable and why use \texttt{while} instead of \texttt{if}?}{Q11: What is a condition variable and why use while instead of if?}}\label{q11-what-is-a-condition-variable-and-why-use-while-instead-of-if}

\textbf{Model answer:} Condition variable pairs with a mutex. Thread calls \texttt{wait()} to atomically release mutex and sleep until signaled. After waking, re-acquires mutex. Use \texttt{while} because: (1) spurious wakeups (real POSIX behavior --- thread wakes without explicit signal), (2) multiple waiters --- condition may be true when signaled but false again by the time this thread executes.

\subsection{Q12: Difference between thread pool\textquotesingle s fixed vs cached vs work-stealing executors?}\label{q12-difference-between-thread-pools-fixed-vs-cached-vs-work-stealing-executors}

\setcodelang{}

\begin{verbatim}
Executors.newFixedThreadPool(n):   fixed n threads, unbounded queue
                                   → predictable resource usage, can queue forever
Executors.newCachedThreadPool():   0 to Integer.MAX threads, 60s TTL
                                   → good for short-lived tasks, can explode memory
Executors.newWorkStealingPool():   ForkJoinPool, work-stealing deques per thread
                                   → best for recursive/tree tasks (divide and conquer)
\end{verbatim}

\section{Senior-Level Discussion Points}\label{senior-level-discussion-points}

\subsection{1. Amdahl\textquotesingle s Law --- Theoretical Speedup Limit}\label{1-amdahls-law--theoretical-speedup-limit}

\visualfigure{../figures/01-concurrency/01-amdahls-law.pdf}{Amdahl's Law: maximum speedup is capped by the serial fraction. Even at 95\% parallel, 256 cores yield only ~17x — the serial 5\% dominates.}

\setcodelang{}

\begin{verbatim}
Speedup = 1 / (S + (1-S)/N)

S = serial fraction of program
N = number of processors

If 5% of your code must be serial:
  N=∞ → max speedup = 1/0.05 = 20x (even with infinite cores!)
\end{verbatim}

\textbf{Implication:} Identify and minimize serial bottlenecks before adding more threads/cores.

\subsection{2. Memory Model Sophistication}\label{2-memory-model-sophistication}

Modern CPUs use TSO (Total Store Order --- x86), allowing stores to be reordered relative to other stores from other cores\textquotesingle{} perspective. ARM/POWER allow more reorderings. High-performance code targeting multiple architectures must use explicit memory fences (\texttt{std::atomic\_thread\_fence}), not just volatile.

\subsection{3. Lock-Free Algorithm Design Challenges}\label{3-lock-free-algorithm-design-challenges}

Beyond ABA: lock-free algorithms are notoriously difficult. Even with CAS, you need careful thought about:

\begin{itemize}
\tightlist
\item
  Memory reclamation (when is it safe to free a node seen by other threads?)
\item
  Progress guarantees (lock-free vs wait-free)
\item
  Platform memory models
\end{itemize}

Epoch-based reclamation and hazard pointers are standard techniques for safe memory reclamation in lock-free data structures.

\subsection{4. Virtual Threads (Java 21 Project Loom)}\label{4-virtual-threads-java-21-project-loom}

Java 21 introduces virtual threads --- millions of lightweight threads mapped to OS threads by the JVM. Blocking operations (I/O, sleep) on virtual threads park the virtual thread without blocking the underlying OS thread. This combines the programming simplicity of threads with async scalability.

\setcodelang{Java}

\begin{Shaded}
\begin{Highlighting}[]
\CommentTok{// Before Loom: thread{-}per{-}request limited to \textasciitilde{}10K concurrent requests}
\CommentTok{// With Loom: 1M virtual threads, each blocking synchronously}
\ControlFlowTok{try} \OperatorTok{(}\DataTypeTok{var}\NormalTok{ executor }\OperatorTok{=} \BuiltInTok{Executors}\OperatorTok{.}\FunctionTok{newVirtualThreadPerTaskExecutor}\OperatorTok{())} \OperatorTok{\{}
    \ControlFlowTok{for} \OperatorTok{(}\DataTypeTok{int}\NormalTok{ i }\OperatorTok{=} \DecValTok{0}\OperatorTok{;}\NormalTok{ i }\OperatorTok{\textless{}} \DecValTok{1\_000\_000}\OperatorTok{;}\NormalTok{ i}\OperatorTok{++)} \OperatorTok{\{}
\NormalTok{        executor}\OperatorTok{.}\FunctionTok{submit}\OperatorTok{(()} \OperatorTok{{-}\textgreater{}} \OperatorTok{\{}
            \BuiltInTok{Thread}\OperatorTok{.}\FunctionTok{sleep}\OperatorTok{(}\DecValTok{1000}\OperatorTok{);}  \CommentTok{// blocks virtual thread, not OS thread}
            \ControlFlowTok{return} \FunctionTok{doWork}\OperatorTok{();}
        \OperatorTok{\});}
    \OperatorTok{\}}
\OperatorTok{\}}
\end{Highlighting}
\end{Shaded}

\subsection{5. STM (Software Transactional Memory)}\label{5-stm-software-transactional-memory}

Database-like transactions for memory: group reads/writes into a transaction, auto-retry on conflict. Clojure uses STM for its reference types. Eliminates deadlocks (transactions can always roll back) but high overhead for write-heavy workloads.

\subsection{6. Reactive/Actor Model}\label{6-reactiveactor-model}

\textbf{Actor model (Akka, Erlang):} Each actor has mailbox and processes one message at a time --- no shared state. Concurrency via message passing. Deadlock-resistant by design (no locks). Actors can be distributed transparently.

\subsection{7. Lock Granularity Trade-offs}\label{7-lock-granularity-trade-offs}

\setcodelang{}

\begin{verbatim}
Coarse-grained locking:    One lock for entire data structure
                           → Simple, but poor concurrency
Fine-grained locking:      Lock per node/segment
                           → Complex, good concurrency, risk of deadlock
Lock-free:                 CAS-based, no locks
                           → Best throughput, very complex to implement correctly
\end{verbatim}

\textbf{ConcurrentHashMap progression:} Per-structure lock \(\rightarrow\) 16 segments \(\rightarrow\) per-bucket CAS/sync (Java 8+).

\section{Typical Mistakes Candidates Make}\label{typical-mistakes-candidates-make}

\subsection{\texorpdfstring{1. Using \texttt{if} instead of \texttt{while} for condition variable waits}{1. Using if instead of while for condition variable waits}}\label{1-using-if-instead-of-while-for-condition-variable-waits}

\textbf{Symptom:} Intermittent NullPointerExceptions or incorrect values after wait(). \textbf{Fix:} Always \texttt{while\ (condition\ not\ met)\ \{\ wait();\ \}}

\subsection{2. Wrong semaphore acquire order causing deadlock}\label{2-wrong-semaphore-acquire-order-causing-deadlock}

\textbf{Symptom:} Producer-consumer deadlocks when buffer is full. \textbf{Fix:} Acquire resource semaphore (empty/full) before mutex, never after.

\subsection{\texorpdfstring{3. Checking for \texttt{==\ null} on volatile reference without synchronized block}{3. Checking for == null on volatile reference without synchronized block}}\label{3-checking-for--null-on-volatile-reference-without-synchronized-block}

\textbf{Symptom:} Two instances created in double-checked locking without volatile. \textbf{Fix:} Declare instance variable as volatile.

\subsection{4. Forgetting to release locks (no RAII/try-finally)}\label{4-forgetting-to-release-locks-no-raiitry-finally}

\textbf{Symptom:} System hangs intermittently when exceptions thrown in critical section. \textbf{Fix:} Always use try-finally or RAII (\texttt{lock\_guard}, \texttt{with\ lock:}).

\subsection{5. Using Thread.stop() or Thread.suspend() (deprecated)}\label{5-using-threadstop-or-threadsuspend-deprecated}

\textbf{Symptom:} Orphaned locks, corrupt state. \textbf{Fix:} Use interrupt flags + cooperative checking (\texttt{while\ (!Thread.interrupted())}).

\subsection{6. Thinking volatile fixes race conditions for compound operations}\label{6-thinking-volatile-fixes-race-conditions-for-compound-operations}

\textbf{Symptom:} Lost updates on \texttt{volatile\ int\ counter} with \texttt{counter++}. \textbf{Fix:} Use \texttt{AtomicInteger.incrementAndGet()}.

\subsection{7. Assuming HashMap.get() is thread-safe because it\textquotesingle s "read-only"}\label{7-assuming-hashmapget-is-thread-safe-because-its-read-only}

\textbf{Symptom:} Infinite loops or NPE during concurrent modification (HashMap can infinite loop on get during rehash in Java). \textbf{Fix:} Use \texttt{ConcurrentHashMap}.

\subsection{\texorpdfstring{8. Ignoring lock ordering \(\rightarrow\) deadlock in production}{8. Ignoring lock ordering \textbackslash rightarrow deadlock in production}}\label{8-ignoring-lock-ordering-rightarrow-deadlock-in-production}

\textbf{Symptom:} Rare deadlocks that only happen under load. \textbf{Fix:} Establish and document global lock acquisition order.

\subsection{9. Creating too many threads for CPU-bound work}\label{9-creating-too-many-threads-for-cpu-bound-work}

\textbf{Symptom:} Performance degrades with 100 threads for CPU work on 8-core machine. \textbf{Fix:} Thread count \(\approx\) core count for CPU-bound.

\subsection{10. Assuming thread-safe classes compose safely}\label{10-assuming-thread-safe-classes-compose-safely}

\textbf{Symptom:} \texttt{if\ (!list.contains(x))\ list.add(x)} on a \texttt{CopyOnWriteArrayList} still has race. \textbf{Fix:} The compound check-then-act still needs external synchronization.

\section{How This Connects To Other Topics}\label{how-this-connects-to-other-topics}

\subsection{Operating Systems}\label{operating-systems}

\begin{itemize}
\tightlist
\item
  Thread scheduling: preemptive vs cooperative, priority inversion
\item
  Kernel mutexes vs userspace futexes (Linux \texttt{futex} syscall --- fast path is CAS in userspace, slow path escalates to kernel)
\item
  Context switching cost \(\rightarrow\) design trade-offs for thread pool sizing
\item
  Virtual memory and copy-on-write (fork() performance)
\end{itemize}

\subsection{System Design}\label{system-design}

\begin{itemize}
\tightlist
\item
  Stateless services avoid shared mutable state --- scale horizontally without concurrency bugs
\item
  Connection pools (JDBC, Redis clients) use semaphores for resource limiting
\item
  Message queues (Kafka, SQS) implement producer-consumer at distributed scale
\item
  Distributed locks (Redis \texttt{SET\ NX\ EX}, ZooKeeper) = distributed mutex
\item
  Optimistic locking (version numbers in DB) = CAS at database level
\item
  MVCC (Multi-Version Concurrency Control) = database reader-writer problem solution
\end{itemize}

\subsection{Performance Engineering}\label{performance-engineering}

\begin{itemize}
\tightlist
\item
  False sharing \(\rightarrow\) memory-conscious data layout (LMAX Disruptor pattern)
\item
  Lock contention profiling (Java VisualVM, Linux \texttt{perf\ lock})
\item
  Amdahl\textquotesingle s Law limits scaling benefits
\item
  Cache-coherent NUMA architectures --- prefer thread-local data
\end{itemize}

\subsection{Databases}\label{databases}

\begin{itemize}
\tightlist
\item
  ACID transactions: serializable isolation = "writer-writer exclusion + reader-writer consistency"
\item
  Pessimistic locking (\texttt{SELECT\ FOR\ UPDATE}) = mutex
\item
  Optimistic locking (version field) = CAS
\item
  Deadlock detection via wait-for graph --- databases kill the youngest transaction
\item
  MVCC: reads don\textquotesingle t block writes, writes don\textquotesingle t block reads --- read-committed, snapshot isolation
\end{itemize}

\subsection{Distributed Systems}\label{distributed-systems}

\begin{itemize}
\tightlist
\item
  CAP theorem: consistency requires cross-node synchronization (distributed mutex)
\item
  Consensus algorithms (Raft, Paxos) = distributed condition variable for leader election
\item
  Distributed transactions (2PC) \(\rightarrow\) distributed deadlock possible
\item
  Event sourcing / CQRS \(\rightarrow\) separates reads and writes to avoid reader-writer contention
\end{itemize}

\section{FAANG Interview Tips}\label{faang-interview-tips}

\subsection{Strategy by Interview Type}\label{strategy-by-interview-type}

\textbf{Coding round:} If given a multi-threading problem:

\begin{enumerate}
\def\labelenumi{\arabic{enumi}.}
\tightlist
\item
  First identify shared state
\item
  Show you understand the race condition before writing fix
\item
  Start with correct-but-simple (mutex), then optimize if asked
\item
  Always mention: "I\textquotesingle d use \texttt{while} not \texttt{if} for wait loops"
\end{enumerate}

\textbf{System design round:}

\begin{itemize}
\tightlist
\item
  Proactively mention concurrency when designing: "The rating service has a counter --- I\textquotesingle d use atomic increment to avoid locks"
\item
  Show awareness of pool sizing: "For I/O-heavy service, I\textquotesingle d size thread pool larger than CPU cores"
\item
  Discuss connection pool limits, semaphores for rate limiting
\end{itemize}

\textbf{Behavioral/technical deep-dive:}

\begin{itemize}
\tightlist
\item
  Have a story about debugging a race condition or deadlock (even in side projects)
\item
  Mention specific tools: Java thread dumps, \texttt{jstack}, Go\textquotesingle s race detector (\texttt{go\ test\ -race}), Helgrind
\end{itemize}

\subsection{Language-Specific Tips}\label{language-specific-tips}

\textbf{Java:} Know \texttt{synchronized}, \texttt{volatile}, \texttt{java.util.concurrent.*}, \texttt{ReentrantLock}, \texttt{ReadWriteLock}, \texttt{AtomicInteger}, \texttt{CountDownLatch}, \texttt{CyclicBarrier}, \texttt{Semaphore}, \texttt{CompletableFuture}. Virtual threads (Java 21).

\textbf{Python:} Know the GIL, \texttt{threading.Lock()}, \texttt{threading.Condition()}, \texttt{multiprocessing} for CPU parallelism, \texttt{asyncio} for I/O concurrency, \texttt{concurrent.futures}.

\textbf{C++:} \texttt{std::mutex}, \texttt{std::lock\_guard}, \texttt{std::unique\_lock}, \texttt{std::condition\_variable}, \texttt{std::atomic}, \texttt{std::memory\_order\_*}, \texttt{std::thread}.

\textbf{Go:} Goroutines, channels (preferred over shared memory), \texttt{sync.Mutex}, \texttt{sync.RWMutex}, \texttt{sync.WaitGroup}, \texttt{sync/atomic}, \texttt{go\ test\ -race}.

\subsection{Common FAANG Patterns They Test}\label{common-faang-patterns-they-test}

\begin{enumerate}
\def\labelenumi{\arabic{enumi}.}
\tightlist
\item
  \textbf{Thread-safe singleton} (almost every Java interview)
\item
  \textbf{Producer-consumer with bounded buffer} (Amazon, Google favorites)
\item
  \textbf{Rate limiter using semaphore} (Uber, Meta)
\item
  \textbf{ReadWriteLock for cache} (high-read, rare-write)
\item
  \textbf{Parallel merge sort / parallel tree traversal} (divide and conquer + thread pool)
\item
  \textbf{Deadlock diagnosis from thread dump} (senior-level Google SRE/SWE)
\end{enumerate}

\subsection{What Separates Good from Great Answers}\label{what-separates-good-from-great-answers}

\begin{longtable}[]{@{}
  >{\raggedright\arraybackslash}p{(\columnwidth - 2\tabcolsep) * \real{0.3316}}
  >{\raggedright\arraybackslash}p{(\columnwidth - 2\tabcolsep) * \real{0.6684}}@{}}
\toprule\noalign{}
\rowcolor{acc!16}
\begin{minipage}[b]{\linewidth}\raggedright
Good Answer
\end{minipage} & \begin{minipage}[b]{\linewidth}\raggedright
Great Answer
\end{minipage} \\
\midrule\noalign{}
\endhead
\bottomrule\noalign{}
\endlastfoot
"Use a mutex" & "Mutex here, but if reads dominate, ReadWriteLock improves throughput" \\
\rowcolor{zebra}
"volatile makes it thread-safe" & "volatile ensures visibility and ordering but not compound atomicity" \\
"Threads run concurrently" & "On N cores, up to N threads run in parallel; others context-switch" \\
\rowcolor{zebra}
"Use thread pool" & "I/O-bound \(\rightarrow\) size at cores\(\times\)(1+wait/compute); CPU-bound \(\rightarrow\) cores+1" \\
"Deadlock: two threads waiting" & "Four Coffman conditions; break circular wait via lock ordering" \\
\end{longtable}

\section{Revision Cheat Sheet}\label{revision-cheat-sheet}

\subsection{10-Minute Revision Summary}\label{10-minute-revision-summary}

\textbf{The 3 fundamental problems:}

\begin{enumerate}
\def\labelenumi{\arabic{enumi}.}
\tightlist
\item
  \textbf{Race conditions} \(\rightarrow\) shared mutable state + no sync \(\rightarrow\) use locks or atomics
\item
  \textbf{Deadlocks} \(\rightarrow\) MHNC conditions \(\rightarrow\) break circular wait with lock ordering
\item
  \textbf{Visibility} \(\rightarrow\) CPU caches \(\rightarrow\) volatile for flags, synchronized/atomic for compounds
\end{enumerate}

\textbf{Key primitives:}

\begin{itemize}
\tightlist
\item
  \textbf{Mutex} = binary lock with ownership, one thread at a time
\item
  \textbf{Semaphore} = counter lock, no ownership, N threads at a time
\item
  \textbf{Condition variable} = wait for condition (ALWAYS while-loop, not if)
\item
  \textbf{volatile} = visibility + ordering, NOT atomicity for compounds
\item
  \textbf{CAS/Atomic} = lock-free compare-and-swap, handles ABA with stamp
\item
  \textbf{SpinLock} = CPU-burning wait, use only for \textless{} 1µs critical sections
\end{itemize}

\textbf{Key algorithms:}

\begin{itemize}
\tightlist
\item
  \textbf{Producer-Consumer:} empty/full semaphores + mutex; acquire resource semaphore BEFORE mutex
\item
  \textbf{Reader-Writer:} ReadWriteLock; readers can share; writer is exclusive
\item
  \textbf{Thread Pool:} queue + N workers; size = cores+1 (CPU) or cores\(\times\)10 (I/O)
\end{itemize}

\textbf{Language facts:}

\begin{itemize}
\tightlist
\item
  \textbf{Python:} GIL = only 1 thread executes Python bytecode; use multiprocessing for CPU parallelism
\item
  \textbf{Java:} \texttt{synchronized} = monitor (mutex + condition); \texttt{volatile} = visibility; \texttt{Atomic*} = CAS
\item
  \textbf{Go:} prefer channels over shared memory; goroutines are cheap (\textasciitilde2KB stack)
\end{itemize}

\subsection{Key Points to Remember}\label{key-points-to-remember}

\begin{itemize}
\tightlist
\item
  Concurrency = dealing with many things at once (design); Parallelism = doing many things at once (execution)
\item
  Thread creation cost: µs; Context switch cost: 100ns (thread) vs µs (process)
\item
  Deadlock 4 conditions: \textbf{MHNC} --- break ANY ONE to prevent
\item
  Lock ordering is the most practical deadlock prevention in software
\item
  Always \texttt{while}, never \texttt{if}, when checking wait condition
\item
  Acquire resource semaphore BEFORE mutex in producer-consumer
\item
  volatile alone is insufficient for \texttt{count++} --- use AtomicInteger
\item
  False sharing: pad to 64 bytes to separate hot variables into own cache lines
\item
  Pool sizing: CPU-bound \(\rightarrow\) N+1; I/O-bound \(\rightarrow\) N\(\times\)(1 + wait/compute)
\item
  Python GIL: threads for I/O, multiprocessing for CPU
\item
  Livelock: both active, no progress; Starvation: one thread never gets resource
\end{itemize}

\subsection{Compact Cheat Sheet Table}\label{compact-cheat-sheet-table}

\begin{longtable}[]{@{}
  >{\raggedright\arraybackslash}p{(\columnwidth - 6\tabcolsep) * \real{0.1323}}
  >{\raggedright\arraybackslash}p{(\columnwidth - 6\tabcolsep) * \real{0.3149}}
  >{\raggedright\arraybackslash}p{(\columnwidth - 6\tabcolsep) * \real{0.2647}}
  >{\raggedright\arraybackslash}p{(\columnwidth - 6\tabcolsep) * \real{0.2880}}@{}}
\toprule\noalign{}
\rowcolor{acc!16}
\begin{minipage}[b]{\linewidth}\raggedright
Concept
\end{minipage} & \begin{minipage}[b]{\linewidth}\raggedright
One-line Definition
\end{minipage} & \begin{minipage}[b]{\linewidth}\raggedright
Key Code / Formula
\end{minipage} & \begin{minipage}[b]{\linewidth}\raggedright
Gotcha
\end{minipage} \\
\midrule\noalign{}
\endhead
\bottomrule\noalign{}
\endlastfoot
Thread & Lightweight execution unit, shared memory & \texttt{new\ Thread(runnable).start()} & Stack is separate; heap is shared \\
\rowcolor{zebra}
Process & Isolated address space & \texttt{fork()} / \texttt{ProcessBuilder} & IPC needed for communication \\
Mutex & Binary lock with ownership & \texttt{lock.acquire();\ ...\ lock.release()} & Forgetting release = deadlock \\
\rowcolor{zebra}
Semaphore & Counter lock, any thread signals & \texttt{sem.wait();\ ...\ sem.signal()} & Acquire resource sem before mutex \\
Condition Var & Sleep + release lock atomically & \texttt{while(!cond)\ cv.wait(lock)} & Use while, not if \\
\rowcolor{zebra}
volatile & Visibility + ordering, no atomicity & \texttt{volatile\ boolean\ ready} & Not enough for \texttt{count++} \\
CAS / Atomic & Hardware read-modify-write & \texttt{AtomicInt.compareAndSet(old,\ new)} & ABA problem \\
\rowcolor{zebra}
Spinlock & Busy-wait loop & \texttt{while(lock.test\_and\_set())\ \{\}} & Wastes CPU; only for short waits \\
Deadlock & Circular wait for resources & 4 conditions: MHNC & Prevent: global lock ordering \\
\rowcolor{zebra}
Livelock & Active but no progress & Randomized backoff & Looks like progress, isn\textquotesingle t \\
Starvation & One thread perpetually denied & Fair queue (FIFO lock) & Different from deadlock \\
\rowcolor{zebra}
Thread Pool & Pre-created workers + task queue & cores+1 (CPU), cores\(\times\)10 (I/O) & Unbounded queue = memory leak \\
False Sharing & Different vars, same cache line & Pad to 64 bytes & Hard to detect; use perf tools \\
\rowcolor{zebra}
Python GIL & One thread at a time in CPython & Use multiprocessing for CPU work & I/O threads still help \\
ReadWriteLock & Multiple readers OR one writer & \texttt{rwLock.readLock()} & Not worth it for write-heavy \\
\rowcolor{zebra}
Producer-Consumer & Bounded buffer sync problem & empty + full semaphores + mutex & Order: resource sem \(\rightarrow\) mutex \\
async/await & Cooperative coroutines on event loop & \texttt{await\ asyncio.gather(*tasks)} & Sync code blocks entire loop \\
\end{longtable}

\subsection{Most Important Interview Concepts (Ranked by Frequency)}\label{most-important-interview-concepts-ranked-by-frequency}

\begin{enumerate}
\def\labelenumi{\arabic{enumi}.}
\tightlist
\item
  \textbf{Race condition definition + example + fix} --- asked everywhere
\item
  \textbf{Deadlock: 4 conditions (MHNC) + prevention} --- almost guaranteed at senior level
\item
  \textbf{Producer-Consumer implementation} --- Amazon, Google standard question
\item
  \textbf{volatile vs synchronized / atomic} --- Java-focused companies
\item
  \textbf{Thread pool sizing} --- system design integration
\item
  \textbf{mutex vs semaphore} --- definitional but important
\item
  \textbf{False sharing} --- performance-focused companies (Databricks, Meta)
\item
  \textbf{Condition variable with while-loop} --- shows real understanding
\item
  \textbf{Python GIL} --- any Python shop
\item
  \textbf{CAS and ABA problem} --- senior/staff level lock-free discussions
\end{enumerate}

\emph{End of Concurrency \& Multithreading Study Notes}
\emph{Estimated study time: 3-4 hours for deep read, 10 minutes for cheat sheet revision}

\end{document}
